% Part B fragment. AI-produced by GPT-6 Astra (OpenAI, ChatGPT Pro mode).
% This fragment contains no preamble. Bibliography keys are in references_B.bib.
\section{Part B: the Picard group scheme of a birational model}

We use $\operatorname{Pic}^{\tau}_{Z/k}$ for the union of the connected
components whose classes in the N\'eron--Severi group are torsion.  All
models in this part are integral, separated, and of finite type over the
indicated field.  A proper model of $K/k$ is supplied with a specified
identification of its function field with $K$.  The ordering of models
is domination by proper birational morphisms; pullback gives the
corresponding direct system of Picard group schemes.  We make no claim
of priority for any argument below.

\subsection{The manuscript's infinitesimal gap and its repair}

On cote 29, scans 071--074, the manuscript reduces the proposed
isomorphism to an infinitesimal question.  The displayed identification
on scan 071 must read $\operatorname{Pic}(Y)\simeq
\operatorname{Pic}(U)$: this is the identification supplied by regularity
of $Y$, and it is the one used in the immediately following diagram.
We have compared the two supplied transcriptions of these pages.
The claimed implication on scan 072 from a bijective closed immersion
and an isomorphism on tangent spaces to an isomorphism of schemes is
false.  This invalid implication does not make the proposed Picard
theorem false.

\subsubsection*{Statement B.1a (Refuted implication; explicit replacement)}
Let $k$ have characteristic $p>0$.  The inclusion
\[
 \alpha_p=\operatorname{Spec}k[t]/(t^p)
 \ \longhookrightarrow\ 
 \alpha_{p^2}=\operatorname{Spec}k[t]/(t^{p^2})
\]
is a closed immersion of commutative $k$-group schemes, is bijective on
geometric points, and induces an isomorphism on tangent spaces at the
identity.  It is not an isomorphism.  Therefore a proof of B.1 must
control more than geometric points and first order tangent vectors.
\begin{proof}
Both groups have the Hopf algebra law $\Delta(t)=t\otimes1+1\otimes t$.
The ideals $(t^{p})$ and $(t^{p^2})$ are Hopf ideals because all the
intermediate binomial coefficients for a $p$-power vanish modulo $p$.
The coordinate homomorphism is the proper quotient
$k[t]/(t^{p^2})\twoheadrightarrow k[t]/(t^p)$.  Both spectra have a
single geometric point.  In both cases the cotangent space of the
augmentation ideal is generated freely by the class of $t$, and the
quotient map is the identity on that class.  The coordinate dimensions
are $p^2$ and $p$, respectively, so the quotient is not an isomorphism.
The accompanying finite certificate checks these assertions for
$p=2,3,5$; the argument just given proves them for every prime.
\end{proof}

\subsubsection*{Statement B.1b (Known purity theorem, with its application)}
Suppose that $f:X\to Y$ is proper and birational, $X$ is normal and
integral, and $Y$ is regular and integral.  For every finite flat
$Y$-group scheme $G$, pullback is an equivalence
\[
 f^*: \operatorname{Tors}_{\mathrm{fppf}}(Y,G)
 \xrightarrow{\ \sim\ }
 \operatorname{Tors}_{\mathrm{fppf}}(X,f^*G).
 \tag{B.1.1}
\]
The groups need not be commutative.  The cited input is Moret--Bailly's
purity theorem, in the detailed proof of
\cite[Theorem 6, pp.~433--434]{Marrama2026}, together with full
faithfulness on normal schemes, Proposition 3 of the same paper.
\begin{proof}
There is an open immersion $j:U\hookrightarrow Y$ with
$\operatorname{codim}_Y(Y\setminus U)\geq2$ such that
$f^{-1}(U)\to U$ is an isomorphism.  To see the codimension assertion,
the inverse rational map extends at every codimension one point of
$Y$ by properness and the valuative criterion, and the two inverse
maps then agree on an open neighbourhood.  Equivalently, we use the
usual codimension one isomorphism locus of a proper birational
morphism to a normal scheme.

Let $P\to X$ be an $f^*G$-torsor.  Its restriction to $f^{-1}(U)=U$
extends to a $G$-torsor $Q\to Y$ by the stated purity theorem.  The
torsors $f^*Q$ and $P$ agree over a dense open of the normal scheme
$X$, hence agree uniquely over $X$ by the cited full faithfulness.
The same full faithfulness, first on $U$ and then on $Y$, shows that
every morphism between pullbacks comes from a unique morphism on
$Y$.  This proves both essential surjectivity and full faithfulness.
\end{proof}

\subsubsection*{Statement B.1c (Proved: finite subgroups detect a proper group)}
Let $k$ be algebraically closed of characteristic $p>0$, let $P$ be a
proper commutative $k$-group scheme of finite type, and let
$i:Q\hookrightarrow P$ be a closed subgroup scheme.  If
\[
 \operatorname{Hom}_{k\text{-groups}}(H,Q)
 \longrightarrow \operatorname{Hom}_{k\text{-groups}}(H,P)
 \tag{B.1.2}
\]
is surjective for every finite commutative $k$-group scheme $H$, then
$i$ is an isomorphism.
\begin{proof}
Since $k$ is perfect, $P_{\mathrm{red}}$ is a subgroup and its identity
component $A$ is an abelian variety.  We use the torsion and divisibility
facts in \cite[Proposition 39.9.11, Tag 03RP]{StacksB}.
For any prime $\ell\ne p$, the
$\ell$-power torsion of $A$ is Zariski dense in $A$.  Indeed, the closure
is a closed subgroup; if its identity component had dimension $b<
\dim A=a$, its finitely many components would contain at most a fixed
multiple of $\ell^{2rb}$ points killed by $\ell^r$, whereas
$A[\ell^r](k)$ has $\ell^{2ra}$ points.  Applying (B.1.2) to these
finite subgroup schemes therefore gives $A\subset Q$.

Every component of $P_{\mathrm{red}}$ contains a torsion point.  For
if $x$ is a point of a component of order $n$ in the finite component
group, then $nx\in A(k)$.  The surjectivity of $[n]:A\to A$ gives
$a\in A(k)$ with $na=nx$, so $x-a$ is torsion and lies in the same
component.  Its finite cyclic subgroup lies in $Q$ by (B.1.2).
Translation by these points shows $P_{\mathrm{red}}\subset Q$.

Choose $r$ so large that the $p^r$-th powers of all nilpotents of
$P$ vanish.  Relative Frobenius
$F^r_{P/k}:P\to P^{(p^r)}$ then factors through
$P_{\mathrm{red}}^{(p^r)}$.  Its kernel $H$ is finite, because
Frobenius on a finite type scheme over a perfect field is finite.
It is commutative and lies in $Q$ by (B.1.2).  Frobenius on the
smooth group $P_{\mathrm{red}}$ is faithfully flat and surjective.
Consequently, for a $T$-point $x$ of $P$, an fppf cover $T'\to T$
and $a\in P_{\mathrm{red}}(T')$ exist with
$F^r(a)=F^r(x)$.  We then have $x-a\in H(T')$, and both summands
belong to $Q(T')$.  Thus $Q\to P$ is surjective as an fppf sheaf.
It is already a monomorphism, so it is an isomorphism.
\end{proof}

\subsubsection*{Statement B.1 (Proved, including every margin variant)}
Let $k$ be algebraically closed of characteristic $p>0$, and let
$f:X\to Y$ be a dominant birational morphism of integral proper
$k$-schemes, with $X$ normal and $Y$ regular.  Then
\[
 f^*: \operatorname{Pic}^{\tau}_{Y/k}
       \xrightarrow{\ \sim\ }
       \operatorname{Pic}^{\tau}_{X/k}
 \tag{B.1.3}
\]
is an isomorphism of group schemes.  Projectivity of $Y$ is
unnecessary.  Over the given algebraically closed field, normality
of $X$ is already geometric normality, and regularity of $Y$ is
equivalent to smoothness over $k$; hence the two other margin
variants have exactly the same conclusion.
\begin{proof}
We write $Q=\operatorname{Pic}^{\tau}_{Y/k}$ and
$P=\operatorname{Pic}^{\tau}_{X/k}$.  Since $f$ is proper birational
and $Y$ is normal, $f_*\mathcal O_X=\mathcal O_Y$.  This equality
persists after every extension $T\to\operatorname{Spec}k$ by flat
base change.  If an invertible sheaf on $Y_T$ becomes trivial after
pullback, the projection formula and this equality show that it
was trivial already; the same assertion after an fppf localization
on $T$ proves that $f^*$ is a monomorphism of relative Picard
sheaves.  Normal proper varieties have proper identity Picard
component, and their torsion component is of finite type
\cite[Theorem 5.4, Remark 5.6, and Remark 6.19]{Kleiman2005}.
Thus $P$ and $Q$ are proper: each has finitely many components, all
translates of its proper identity component.  A proper
monomorphism is a closed immersion, so $f^*:Q\hookrightarrow P$
is a closed immersion.

For a finite commutative $k$-group scheme $H$, write
$H^D=\underline{\operatorname{Hom}}(H,\mathbb G_m)$ for its
Cartier dual.  Raynaud's Picard--torsor duality gives the natural
isomorphism
\[
 H^1_{\mathrm{fppf}}(Z,H^D)
 \simeq
 \operatorname{Hom}_{k\text{-groups}}(H,\operatorname{Pic}_{Z/k})
 \tag{B.1.4}
\]
for $Z=X,Y$; see \cite[Proposition 6.2.1, pp.~50--51]{Raynaud1970}.
Here $Z\to\operatorname{Spec}k$ is proper, flat and finitely
presented, with $H^0(Z,\mathcal O_Z)=k$ universally.  A $k$-point
of $Z$ and the fact that every finite torsor over an algebraically
closed field has a point identify the relative $R^1$ in Raynaud's
statement with the displayed absolute $H^1$.
Every group homomorphism from a finite group to
$\operatorname{Pic}_{Z/k}$ lands in its torsion component, because
its image in the component group is finite.  Statement B.1b,
applied to $G=H^D\times_kY$, therefore makes (B.1.2) bijective
for $Q\hookrightarrow P$.  Statement B.1c completes the proof.
This argument controls arbitrary finite subgroup schemes,
including nonreduced ones, so it supplies the missing
infinitesimal step of the manuscript.
\end{proof}

\subsubsection*{Statement B.1d (Proved: the manuscript's tangent-space calculation)}
Under the hypotheses of B.1, pullback
$H^1(Y,\mathcal O_Y)\to H^1(X,\mathcal O_X)$ is an isomorphism.
The Frobenius argument on scans 073--074 does establish this
claim, independently of the invalid scheme-theoretic inference
on scan 072.
\begin{proof}
The map is injective because $f_*\mathcal O_X=\mathcal O_Y$.
Let $U\subset Y$ be the codimension one isomorphism locus from
B.1b.  For $n\geq1$ and $Z=X,Y$, set
$B_{n,Z}=F^n_*\mathcal O_Z/\mathcal O_Z$.  This coherent sheaf
is torsion free if $Z$ is normal: a rational $p^n$-th root
of a regular function is integral, hence regular.  If $Z=Y$
it is locally free.  Indeed, Frobenius is finite flat on
the regular variety $Y$ by Kunz's theorem
\cite[Tag 0EC0]{StacksB}; in the resulting finite free
module the element $1$ is unimodular, since it does not
belong to the maximal ideal times that module.  The quotient
by its rank one summand is therefore locally free.

A global section of $B_{n,X}$ is determined by its restriction
to $f^{-1}(U)$, by torsion freeness.  That restriction extends
uniquely to $B_{n,Y}$ by Hartogs extension.  Conversely the
natural pullback of such a section to $X$ restricts to the
original section on the dense open and hence agrees everywhere.
Thus $H^0(Y,B_{n,Y})\simeq H^0(X,B_{n,X})$ for every $n$.
Since $F^n:k\to k$ is surjective, these spaces identify with
the kernels of $F^n$ on the respective $H^1(\mathcal O)$.
This proves the isomorphism on the Frobenius-nilpotent summands.
For the bijective summands, Artin--Schreier identifies a basis
of Frobenius-fixed vectors with $H^1_{\text{ét}}
(-,\mathbb F_p)$, and B.1b for the constant group
$\mathbb F_p$ identifies those groups.  The semilinear
decomposition gives the desired full isomorphism.
\end{proof}

\subsubsection*{Statement B.1e (Proved): the regular target over an imperfect field}
Let $k$ be any field of characteristic $p>0$.  Let $f:X\to Y$ be a
proper birational morphism of proper geometrically integral $k$-schemes,
with $X$ geometrically normal and $Y$ regular.  Then pullback is an
isomorphism on $\operatorname{Pic}^{\tau}$.
\begin{proof}
The equality $f_*\mathcal O_X=\mathcal O_Y$ commutes with field extension.
Since $X_{\bar k}$ is normal, this equality shows that $Y_{\bar k}$ is
normal as well: a rational function integral over an affine open of
$Y_{\bar k}$ pulls back to a regular function on its inverse image in
$X_{\bar k}$ and then descends by the displayed equality.  Thus both
Picard torsion components are proper and of finite type, and pullback
is a closed immersion $Q\hookrightarrow P$ as in B.1.

The smooth locus of $X/k$ is nonempty, so it has a point over some
finite separable extension $l/k$.  Normality of $X$ and regularity of
$Y$ survive that extension.  It suffices to prove the assertion over
$l$, and we therefore assume that $X(k)\ne\varnothing$; its image gives
a point of $Y$ as well.  For every finite commutative $k$-group $H$,
Raynaud's Proposition 6.2.1 and the sections of the structure maps give
\[
 \operatorname{Hom}(H,\operatorname{Pic}_{Z/k})
 \simeq H^1(Z,H^D)/H^1(k,H^D),\qquad Z=X,Y.
\]
Indeed, the Leray edge map is surjective because the section makes
$H^2(k,H^D)\to H^2(Z,H^D)$ injective.  B.1b therefore implies that
$\operatorname{Hom}(H,Q)\to\operatorname{Hom}(H,P)$ is bijective.

We can detect the closed subgroup over $k$ itself.  For each positive
integer $n$, the group $P[n]$ is finite: after an algebraic closure
this follows from finiteness of multiplication on the reduced abelian
identity component and the finitely many components.  Hence $P[n]$
lies in $Q$.  After extending to $\bar k$, the torsion-density and
torsion-representative argument of B.1c gives
$(P_{\bar k})_{\mathrm{red}}\subset Q_{\bar k}$.
The finite kernel of every relative Frobenius power of $P$ also lies
in $Q$.  Choose a power killing the nilradical of $P_{\bar k}$ and use
the final fppf argument of B.1c.  It gives $Q_{\bar k}=P_{\bar k}$,
hence $Q=P$.  We never apply normality of $X$ after a purely
inseparable extension in the torsor-purity argument.
\end{proof}


\subsection{A uniform curve and a uniform bound for first cohomology}

\subsubsection*{Statement B.2a (Proved: a geometric generic curve detects forms)}
Let $k$ be algebraically closed of characteristic $p>0$, and let
$M$ be a normal projective integral $k$-scheme of dimension $d\geq2$.
Let $U=M_{\mathrm{reg}}$ and
$V=H^0(U,\Omega^1_{U/k})$.  There are a field extension $E/k$ with
$E$ algebraically closed and a smooth proper connected curve
$C\hookrightarrow U_E$ with these properties:
\begin{enumerate}
\item $V\otimes_kE\to H^0(C,\Omega^1_{C/E})$ is injective.
\item For every closed $k$-subscheme $Z\subset M$ of codimension at
least two, $C\cap Z_E=\varnothing$.
\item For every finite surjection $q:N\to M$ with $N$ integral,
the inverse image $q_E^{-1}(C)$ is geometrically irreducible.
\end{enumerate}
The curve is the geometric generic complete intersection of
$d-1$ hypersurfaces of one fixed sufficiently large degree.
The quantifiers in (2) and (3) refer to objects defined over $k$;
we do not assert these properties for all objects defined over $E$.
\begin{proof}
Choose a very ample invertible sheaf $\mathcal O_M(1)$.  Normality
gives a coherent reflexive sheaf
$\mathcal F=j_*\Omega^1_{U/k}$, where $j:U\hookrightarrow M$,
and $V=H^0(M,\mathcal F)$ is finite dimensional.  Choose $m$ so
that $\mathcal F^\vee(m)$ is generated by finitely many global
sections.  Dualizing the resulting surjection and using
$\mathcal F=\mathcal F^{\vee\vee}$ gives
\[
 \mathcal F\hookrightarrow\mathcal O_M(m)^{\oplus r}.
 \tag{B.2.1}
\]
On $U$ this is a locally split injection of vector bundles.

Fix $n>m$ for which $\mathcal O_M(n)$ is very ample.  Let $L$ be
the function field of the product of $d-1$ copies of the
parameter space of hypersurfaces in
$H^0(M,\mathcal O_M(n))$.  Successively cutting by the universal
generic members gives $M_L=M_0\supset M_1\supset\cdots
\supset M_{d-1}=C_L$.  Each $M_i$ is integral over the field
containing the parameters introduced through stage $i$, and
remains integral after the further purely transcendental
extensions of parameters.  Indeed, for an integral scheme with
a base point free linear system the universal incidence divisor
is a projective bundle over that scheme; it is integral, and its
generic fibre over the parameter space is integral.  This proves
the assertion at each stage without requiring normality of the
intermediate hyperplane sections.

At each stage, the divisor sequence, twisted by $m$, has the form
\[
 0\longrightarrow\mathcal O_{M_i}(m-n)
 \longrightarrow\mathcal O_{M_i}(m)
 \longrightarrow\mathcal O_{M_{i+1}}(m)
 \longrightarrow0.
\]
The first term has no global sections: it is the negative of a
positive power of an ample sheaf on an integral proper scheme
of positive dimension.  Thus
$H^0(M,\mathcal O_M(m))\otimes_kL\to
H^0(C_L,\mathcal O_{C_L}(m))$ is injective.  By (B.2.1),
restriction of $V\otimes_kL$ to
$H^0(C_L,\Omega^1_{U_L/L}|_{C_L})$ is injective too.

Bertini smoothness and geometric irreducibility apply to these
generic members: we use \cite[Tags 0FD6 and 0G4F]{StacksB}.
The codimension of $M\setminus U$ is at least two, so a generic
complete intersection of $d-1$ hypersurfaces misses it.  After
extending $L$ to an algebraic closure $E$, $C=C_L\times_LE$ is
therefore smooth, proper, and geometrically integral.  Flat base
change preserves the injection just proved.  Its conormal sequence
in $U_E$ is
\[
 0\longrightarrow\mathcal O_C(-n)^{\oplus(d-1)}
 \longrightarrow\Omega^1_{U_E/E}|_C
 \longrightarrow\Omega^1_{C/E}\longrightarrow0.
\]
The first sheaf has no global sections.  This proves (1).

For a fixed closed $Z\subset M$ of dimension at most $d-2$,
successive general hypersurfaces reduce its dimension until the
intersection is empty.  The locus of such tuples is a nonempty
open subset of the parameter space, hence contains its generic
point.  This proves (2) for every $k$-defined $Z$, without
intersecting a collection of open subsets at a closed point.

For (3), we apply the geometric irreducibility form of Bertini
to the pullback linear system on $N$, and then successively to
its generic sections.  Since $q$ is finite, this linear system
is base point free and its image has the same dimension as the
scheme being cut.  At every stage that dimension is at least two.
The precise hypotheses of Stacks Tag 0G4F are satisfied: two
general members have a component of codimension two, and a third
member misses its generic point; the base locus is empty.
The conclusion at each stage is geometric irreducibility, which
does not require reducedness.  Iteration gives (3).
\end{proof}

\subsubsection*{Statement B.2b (Proved: the same curve detects $H^1$ on every modification)}
With $M,E,C$ as in B.2a, let $a:W\to M$ be any proper birational
morphism with $W$ normal and integral.  The inverse of $a$ is
defined on $C$, and the resulting map $i_W:C\to W_E$ induces an
injection
\[
 i_W^*:H^1(W_E,\mathcal O_{W_E})
 \hookrightarrow H^1(C,\mathcal O_C).
 \tag{B.2.2}
\]
\begin{proof}
There is a $k$-defined open $M^\circ\subset M$, containing every
codimension one point, over which $a$ is an isomorphism.  By
B.2a(2), $C\subset M^\circ_E$, so $i_W$ is defined.
Let $T$ be the kernel of (B.2.2).  Absolute Frobenius acts
$p$-semilinearly on both cohomology groups and commutes with
restriction, so $T$ is stable under this action.  Because $E$
is perfect, the finite dimensional semilinear Fitting decomposition
writes $T=T_{\mathrm{nil}}\oplus T_{\mathrm{bij}}$, where Frobenius
is nilpotent on the first summand and bijective on the second.
We show that both summands vanish.

Suppose first that $T_{\mathrm{nil}}\ne0$.  It contains
$0\ne c$ with $F(c)=0$.  For a normal integral $E$-scheme $Z$,
put $B_Z=F_*\mathcal O_Z/\mathcal O_Z$, with the first map
$a\mapsto a^p$.  The derivative gives an injection of this
sheaf into the rational differentials $F_*\Omega^1_{E(Z)/E}$:
if $db=0$ then $b$ is a $p$-th power in $E(Z)$, and its
$p$-th root, being integral over every local ring in which $b$
is regular, is regular by normality.  As
$H^0(W_E,\mathcal O_{W_E})=E$ and $F:E\to E$ is onto, the
cohomology sequence identifies
\[
 H^0(W_E,B_{W_E})
 \simeq \ker\bigl(F:H^1(W_E,\mathcal O_{W_E})
                         \to H^1(W_E,\mathcal O_{W_E})\bigr).
 \tag{B.2.3}
\]
The class $c$ therefore gives a nonzero rational differential
$\omega$.  It is regular on $(M^\circ\cap U)_E$, and extends
across the codimension at least two complement to a section of
$\Omega^1_{U_E/E}$ by Hartogs extension for a vector bundle on
the regular scheme $U_E$.  Flat base change identifies this
space with $V\otimes_kE$.  The identifications in (B.2.3) are
functorial for restriction to $C$.  Since $c|_C=0$, the
restricted differential $\omega|_C$ is zero.  B.2a(1) gives
$\omega=0$, contradicting (B.2.3).  Hence $T_{\mathrm{nil}}=0$.

Suppose next that $T_{\mathrm{bij}}\ne0$.  A bijective
$p$-semilinear endomorphism over an algebraically closed field
has a basis of fixed vectors; in particular there is
$0\ne c\in T$ with $F(c)=c$.  The Artin--Schreier exact sequence
and the surjectivity of $F-1:E\to E$ identify these fixed vectors
with $H^1_{\text{ét}}(W_E,\mathbb F_p)$.  For proper
$W$, extension of algebraically closed ground fields preserves
this latter group; here this follows already from semilinear
linear algebra and flat base change on coherent cohomology
\cite[Lemmas 59.63.3 and 59.63.6, Tag 0A3P]{StacksB}.
Thus $c$ comes from a nontrivial $\mathbb F_p$-torsor $P\to W$
defined over $k$.

This torsor is connected.  Indeed, translations act transitively
on its connected components, and each component is finite
\'etale of a constant positive degree over connected $W$.
The number of components divides $p$; if it were $p$, each
component would have degree one and give a section, making
the torsor trivial.  Since $W$ is normal, $P$ is normal; a
connected normal finite type scheme is integral.  Let $N$ be
the normalization of $M$ in the finite field extension
$k(P)/k(M)$.  It is finite over $M$, since a variety over a
field is excellent.  Over $M^\circ$ it agrees with
$P|_{M^\circ}$, by uniqueness of normalization.  B.2a(3) makes
$N_E\times_{M_E}C$ geometrically irreducible.  But this scheme
is $P_E|_C$, a finite \emph{\'etale} cover of $C$, so it is
reduced and connected.  Consequently that $\mathbb F_p$-torsor
is nontrivial.  This contradicts $c|_C=0$.  Hence
$T_{\mathrm{bij}}=0$ also, and (B.2.2) is injective.
\end{proof}

\subsubsection*{Statement B.2 (Proved for the stated boundedness question)}
For every finitely generated extension $K/k$ of characteristic
$p>0$, there is an integer $b(K/k)$ such that
\[
 \dim_k H^1(X,\mathcal O_X)\leq b(K/k)
 \tag{B.2.4}
\]
for every proper geometrically normal model $X$ of $K/k$.
For algebraically closed $k$ and transcendence degree at least
two, we may take $b(K/k)=g(C)$ for the single curve of B.2a
constructed on any fixed normal projective model of $K$.
\begin{proof}
First suppose $k$ algebraically closed.  If the transcendence
degree is zero, the only model is $\operatorname{Spec}k$ and
the bound is zero.  If it is one, every normal proper model is
the unique smooth proper curve of the function field; its
genus is a bound.  Suppose the transcendence degree is at
least two and fix a normal projective model $M$.
For a given proper normal model $X$, let $W$ be the normalization
of the closure of the graph of the birational map $M\dashrightarrow X$.
There are proper birational morphisms $a:W\to M$ and $b:W\to X$.
Since $X$ is normal, $b_*\mathcal O_W=\mathcal O_X$; the
low degree Leray sequence therefore gives an injection
$H^1(X,\mathcal O_X)\hookrightarrow H^1(W,\mathcal O_W)$.
Flat base change and B.2b then give
\[
 \dim_k H^1(X,\mathcal O_X)
 \leq\dim_E H^1(W_E,\mathcal O_{W_E})
 \leq\dim_E H^1(C,\mathcal O_C)=g(C).
\]

Now let $k$ be arbitrary.  If there is no proper geometrically
normal model the assertion is vacuous.  Otherwise let $k_0$
be the algebraic closure of $k$ in $K$.  On every such model
$X$, normality and properness identify
$H^0(X,\mathcal O_X)$ with $k_0$: an element of $K$ algebraic
over $k$ is integral over each local ring, and hence is regular;
the reverse inclusion follows from properness.  The extension
$k_0/k$ is finite and separable, because geometric normality
implies geometric reducedness.  After extension to an algebraic
closure $\bar k$, $X_{\bar k}$ is a disjoint union of finitely
many normal integral proper schemes indexed by the embeddings
$\sigma:k_0\hookrightarrow\bar k$.  Their function fields
$K_\sigma=\operatorname{Frac}(K\otimes_{k_0,\sigma}\bar k)$
depend only on $K/k$, and each is a regular function field
over $\bar k$.  Apply the algebraically closed case to each
$K_\sigma$ and add the resulting bounds.  Flat base change
and additivity of cohomology on a finite disjoint union give
(B.2.4), with
$b(K/k)=\sum_\sigma b(K_\sigma/\bar k)$.
\end{proof}

\subsubsection*{Statement B.2c (Proved with the indicated scope: tangent stabilization)}
Choose a geometrically normal proper model whose $h^1(\mathcal O)$
is maximal among the bounded dimensions of B.2.  Every geometrically
normal proper model dominating it has the same first cohomology,
via pullback.  If $k$ is perfect, the normal-model category is
filtered, and this is eventual constancy of its tangent system.
Over an arbitrary field, every proper birational morphism from a
geometrically normal model to a smooth proper model induces an
isomorphism on $\operatorname{Pic}^{\tau}$.
\begin{proof}
Pullback on $H^1$ is injective by Leray and $f_*\mathcal O=\mathcal O$.
Maximal dimension therefore makes it an isomorphism.  Over a perfect
field, normalization of a graph is geometrically normal and gives a
common dominating model.  For the last assertion, base change to an
algebraic closure and apply B.1 to each geometric component, then
use faithfully flat descent of an isomorphism.

A tangent bound alone does not imply a group-scheme stabilization
bound: the proper finite groups $\alpha_{p^r}$ form a strictly
increasing system with constant tangent dimension one.  The following
statements supply the additional geometric control over an
algebraically closed field.  Over an imperfect field, a normalized
graph need not be geometrically normal; we do not silently use that
false cofinality assertion.
\end{proof}

\subsubsection*{Statement B.2d (Proved): simultaneous embedding into one Jacobian}
Under the hypotheses and choices of B.2a, write $J=\operatorname{Pic}^0_{C/E}$.
For every proper normal modification $W\to M$, the unique lift
$C\to W_E$ induces a closed immersion
\[
 \operatorname{Pic}^{\tau}_{W/k}\times_kE\hookrightarrow J.
 \tag{B.2.J}
\]
These immersions commute with pullback between modifications.
\begin{proof}
The curve lies in the isomorphism locus of $W\to M$, and this determines
its lift uniquely.  Restriction takes torsion-equivalent line bundles to
degree-zero line bundles, so the displayed map is defined.  Its source
is a proper commutative group of finite type.  Its kernel has zero Lie
algebra by B.2b.  A finite type group scheme over the algebraically closed
field $E$ with zero Lie algebra is reduced of dimension zero: at the
identity its maximal ideal satisfies $\mathfrak m/\mathfrak m^2=0$,
hence $\mathfrak m=0$ by Nakayama's lemma, and translation treats every
closed point.  The kernel is proper, so it is finite \'etale.

If this kernel were nontrivial, its group of $E$-points would have a
nonzero element of prime order $\ell$.  Write
$P=\operatorname{Pic}^{\tau}_{W/k}$.  The group scheme $P[\ell]$ is
finite, even when $\ell=p$: on every component of $P_{\mathrm{red}}$,
multiplication by $\ell$ has finite fibres, because its identity component
is an abelian variety.  Thus $[\ell]:P\to P$ is proper and quasi-finite,
hence finite.  Every $E$-point of the finite $k$-scheme $P[\ell]$ is a
$k$-point, since $k$ is algebraically closed.  Consequently our element
is represented by a line bundle $L$ on $W$ of exact order $\ell$, whose
restriction to $C$ is trivial.

Suppose first that $\ell\ne p$.  The Kummer sequence, together with
$\Gamma(W,\mathcal O_W)^*=k^*$, gives a nontrivial
$\mu_\ell$-torsor $Q\to W$.  It is connected, since a torsor under
a cyclic group of prime order is either connected or trivial.  It is
normal and therefore integral.  Its normalization over $M$ in $k(Q)$
is finite and integral.  By B.2a its inverse image over $C$ is
geometrically irreducible, hence the \'etale cover $Q_E|_C$ is connected.
But $L|_C$ is trivial and $E^*$ is $\ell$-divisible, so this cover is
trivial.  This is a contradiction.

Now let $\ell=p$.  The Kummer sequence gives a nontrivial
$\mu_p$-torsor on $W$.  Restriction to the generic point is injective
for finite torsors on a normal integral scheme, by the normal full
faithfulness used in B.1b.  Its generic class is therefore represented
by $a\in k(W)^*$ with $a\notin k(W)^{*p}$.  The logarithmic differential
$da/a$ is nonzero, since the kernel of $d:k(W)\to
\Omega^1_{k(W)/k}$ is $k(W)^p$.  Local equations for the torsor are
$t^p=u$ with $u$ a unit; changing a trivialization multiplies $u$ by a
$p$-th power.  Hence $da/a$ is regular on $W_{\mathrm{reg}}$.
On the common open of $M$ and $W$ it extends, by Hartogs on the smooth
scheme $U$, to an element of $V=H^0(U,\Omega^1_{U/k})$.
The triviality of $L|_C$, and the surjectivity of $E^*\xrightarrow{p}E^*$,
make the restricted $\mu_p$-torsor trivial.  Thus $da/a$ restricts to
zero on $C$, contradicting B.2a(1).

The kernel is therefore trivial.  A homomorphism of group schemes with
trivial scheme-theoretic kernel is a monomorphism, and a proper
monomorphism is a closed immersion.  Compatibility follows from
uniqueness of the lifted curve.
\end{proof}

\subsubsection*{Statement B.2e (Proved): finite quotients of a smooth model}
Let $k$ be algebraically closed of characteristic $p>0$, and let $Y$ be a smooth projective integral
$k$-scheme, and let a finite group $\Gamma$ act faithfully on $Y$.
Set $K=k(Y)^\Gamma$.  The direct system of
$\operatorname{Pic}^{\tau}$ of proper normal models of $K/k$ is
essentially constant.
\begin{proof}
The quotient $T=Y/\Gamma$ is a normal projective scheme and
$q:Y\to T$ is finite.  We write $n=|\Gamma|$ and
$P=\operatorname{Pic}^{\tau}_{Y/k}$.
For a normal proper modification $W\to T$, let $Y_W$ be the
normalization of $W$ in $k(Y)$.  It is finite over $W$ and has a
proper birational morphism to $Y$; its $\Gamma$-quotient is $W$.
By B.1, pullback identifies $\operatorname{Pic}^{\tau}_{Y_W/k}$
with $P$, compatibly with $\Gamma$ and with further modifications.
Thus pullback gives a map
\[
 h_W:P_W:=\operatorname{Pic}^{\tau}_{W/k}\longrightarrow
 A:=P^\Gamma.
\]
Here $A$ is the closed fixed subgroup scheme.

The kernel of $h_W$ embeds in
$D(\Gamma^{\mathrm{ab}})=\underline{\operatorname{Hom}}(\Gamma,\mathbb G_m)$.
Indeed, after an fppf base change on a test scheme, a line bundle in
the kernel has a trivialization on $Y_W$.  Its canonical
$\Gamma$-linearization becomes a character, because
$\Gamma(Y_W\times_kS,\mathcal O)=\Gamma(S,\mathcal O)$ universally.
Changing the trivialization by a scalar does not change this character.
If it is trivial, the trivializing section and its inverse are invariant;
they descend because $(q_*\mathcal O_{Y_W})^\Gamma=\mathcal O_W$
universally after base change over $k$.  The descended section
trivializes the original line bundle.  This proves injectivity as
fppf sheaves.  The kernel is proper, so this monomorphism into the
finite group $D(\Gamma^{\mathrm{ab}})$ is a closed immersion.
Its order is at most $|\Gamma^{\mathrm{ab}}|\le n$.

Let $B_T\subset A$ be the image of $h_T$.  We claim that
$A_{\mathrm{red}}^0\subset B_T$.  Take $[L]\in A(k)$ and choose
a Cartier divisor $D$ representing $L$ on the regular scheme $Y$.
The divisor $D_\Sigma=\sum_{\gamma\in\Gamma}\gamma D$ is invariant.
At a prime divisor above a prime divisor of $T$, its coefficient is
divisible by the order of the decomposition group, hence by the
ramification index.  Therefore $D_\Sigma=q^*E$ for an integral Weil
divisor $E$ on $T$.  The divisor $nE$ is Cartier.  To verify this at
$x\in T$, trivialize $D_\Sigma$ on the semilocal scheme at the
finitely many points above $x$ and represent it by a rational
function $u$.  After shrinking around $x$, its norm has divisor
$q_*D_\Sigma=nE$.  Thus $nE$ is locally principal.
Since $[L]$ is $\Gamma$-invariant, we have
\[
 q^*\mathcal O_T(nE)\simeq L^{\otimes n^2}.
\]
If $L$ is torsion-equivalent to zero, this pullback is numerically
trivial.  Finite surjective pullback detects numerical triviality,
by degrees on integral curves.  Numerical triviality is equivalent
to torsion equivalence for projective schemes
\cite[Theorem 6.3]{Kleiman2005}.  Hence
$\mathcal O_T(nE)\in P_T(k)$, and $n^2[L]\in B_T(k)$.
Multiplication by $n^2$ is surjective on the abelian variety
$A_{\mathrm{red}}^0$, so the claim follows by density of its $k$-points.

The quotient $A/B_T$ is therefore a fixed finite group scheme, of
some order $N$.  Since $P_T\hookrightarrow P_W$ is a closed immersion,
the quotient $P_W/P_T$ maps to $A/B_T$; its kernel is a quotient
of $\ker(h_W)$.  For clarity, if $h_W(x)$ belongs to $B_T$, then
fppf locally it equals $h_T(y)$ for some $y\in P_T$, because a
homomorphism onto its scheme-theoretic group image is faithfully flat.
Thus $x-y\in\ker h_W$, and the kernel is exactly
$\ker(h_W)/(\ker(h_W)\cap P_T)$.  Quotients by closed normal
subgroups exist as group schemes in this finite-type commutative
setting \cite[Theorems 5.14, 5.28 and 5.39]{MilneAG2022B}.  Hence $P_W/P_T$ is finite of order at most $nN$.
As $W$ varies, these finite quotients form a directed system of
closed immersions with uniformly bounded orders.  Choose one whose
order is maximal.  Every later closed immersion has equal order,
hence is an isomorphism.  This proves eventual constancy.  Normalized
graphs show that modifications of $T$ form a cofinal subsystem of
all proper normal models of $K/k$.
\end{proof}

\subsubsection*{Statement B.2f (Proved): descent through a purely inseparable extension}
Let $k$ be algebraically closed of characteristic $p>0$, let $K'/K$
be a finite purely inseparable extension of finitely generated
$k$-fields, and assume that the Picard system for $K'/k$ is
essentially constant.  Then the same holds for $K/k$.
\begin{proof}
In transcendence degrees zero and one, the unique normal proper model
proves the assertion.  We may therefore assume transcendence degree at
least two.  Choose $e$ with $(K')^{p^e}\subset K$, and put $q=p^e$.
Choose a normal projective model $T'$ of $K'$ after which its
Picard system is constant.  The relative Frobenius twist
$(T')^{(q)}$ has function field identified with $(K')^q$;
normalizing it in the intermediate field $K$ gives a normal
projective model $T$.  The normalization of $T$ in $K'$ is
exactly $T'$.  Thus $\pi:T'\to T$ is finite and purely inseparable.

For any normal proper modification $W\to T$, its normalization
$W'$ in $K'$ is finite over $W$ and maps birationally to $T'$.
The assumed stabilization identifies
$P_{W'}=\operatorname{Pic}^{\tau}_{W'/k}$ with the fixed group
$P'=\operatorname{Pic}^{\tau}_{T'/k}$.  Write
$h_W:P_W\to P'$ for pullback, and $B_T$ for the image of $h_T$.

The field inclusion $(K')^q\subset K$ and normality imply
$\mathcal O_{W'}^q\subset\mathcal O_W$ locally.  They give a
morphism $r_W:W\to(W')^{(q)}$ with
\[
 r_W\pi_W=F^e_{W'/k},\qquad
 \pi_W^{(q)}r_W=F^e_{W/k}.
\]
The two factorizations of relative Frobenius imply, on Picard
group schemes, that $[q]$ kills $\ker h_W$ and that
$[q]P'\subset B_T$.  More explicitly, relative Frobenius on a
scheme induces Verschiebung on its Picard scheme, and composing
this map with relative Frobenius of the Picard scheme is $[q]$;
the displayed factorizations give the two assertions by
functoriality.  Equivalently they follow directly by raising
transition functions of line bundles and their trivializations
to their $q$-th powers.

The second assertion contains $(P')_{\mathrm{red}}^0$ in $B_T$,
since $[q]$ is surjective on that abelian variety.  Therefore
$P'/B_T$ is finite, of some fixed order $N$.
Apply B.2d with $M=T$.  After the one fixed extension $E/k$,
every $P_W$ embeds in the same Jacobian $J$.  Since $[q]$
kills $\ker h_W$, its base change embeds in $J[q]$, whose
order is $q^{2g}$, where $g=\dim J$.  Thus
$|\ker h_W|\le q^{2g}$ uniformly in $W$.

The proper quotient $P_W/P_T$ maps to $P'/B_T$ with kernel a
quotient of $\ker h_W$, so it is finite of order at most
$q^{2g}N$.  The maximal-order argument of B.2e now proves
essential constancy over modifications of $T$, hence over all
normal proper models of $K/k$.
\end{proof}

\subsubsection*{Statement B.2g (Proved): full stabilization over an algebraically closed field}
If $k$ is algebraically closed of positive characteristic and
$K/k$ is finitely generated, the full direct system
$\operatorname{Pic}^{\tau}_{X/k}$ over proper normal models of
$K/k$ is essentially constant.
\begin{proof}
By the regular Galois alteration theorem of de Jong
\cite[Situation 5.3, Theorem 5.13 and Corollary 5.15, pp.~613, 619--620]{deJong1997},
there are a smooth projective integral $Y$ and an action of a
finite group $\Gamma$ on $Y$ such that
$K'=k(Y)^\Gamma$ is finite purely inseparable over $K$.
The construction is projective: its semistable fibrations and
the equivariant modifications in Proposition 5.11 are projective,
starting with a projective model of $K$.
Dividing $\Gamma$ by the kernel of its action does not change
$K'$.  B.2e proves stabilization for $K'/k$, and B.2f descends
it to $K/k$.
\end{proof}

% Checked addition; independently audited by root on 2026-11-20.
\subsubsection*{Statement B.1f (Proved): finite-group detection over an arbitrary field}
Let $k$ be a field of characteristic $p>0$, let $P$ be a commutative
group scheme of finite type over $k$, and let $Q\subset P$ be a closed
subgroup scheme.  If every finite $k$-subgroup scheme of $P$ is contained
in $Q$, then $Q=P$.
\begin{proof}
For each $r\geq 1$, the relative Frobenius $F^r:P\to P^{(p^r)}$ is
finite.  On an affine chart this follows because a finitely generated
$k$-algebra is a finite module over the subalgebra generated by $k$
and all its $p^r$-th powers.  Thus $\ker F^r$ is a finite $k$-group
scheme and is contained in $Q$.

After extending to $\bar k$, put
$R=\mathcal O_{P_{\bar k},e}$ and let $\mathfrak m$ be its maximal
ideal.  The ideal of $\ker F^r$ at $e$ is generated by $p^r$-th powers
of elements of $\mathfrak m$.  The ideal $I$ of $Q_{\bar k}$ therefore
satisfies
\[
 I\subset \mathfrak m^{p^r}\quad\text{for every }r,
 \qquad I\subset\bigcap_{n\geq1}\mathfrak m^n=0.
\]
Here the last equality is the Krull intersection theorem.  Consequently
$Q_{\bar k}$ contains an open neighbourhood of the identity, and is
an open and closed subgroup of $P_{\bar k}$.  In particular it contains
the entire scheme-theoretic identity component.

We show that it meets every geometric component.  Set
$G=(P_{\bar k})_{\mathrm{red}}^0$.  The structure theorem for smooth
commutative algebraic groups gives a smooth connected unipotent subgroup
$U\subset G$ with $G/U$ semiabelian
\cite[Theorem 8.27 and Theorem 16.13]{MilneAG2022B}.
There is an integer $a$ such that $[p^a]U=0$: embed $U$ into a group
of upper triangular unipotent matrices and take $p^a$ at least the
matrix size.  For any $x\in P(\bar k)$, choose $m>0$ such that
$mx\in G(\bar k)$.  Multiplication by $m$ is surjective on the
semiabelian variety $G/U$.  We can therefore choose $y\in G(\bar k)$
with $m(x-y)\in U(\bar k)$.  The point $t=x-y$ lies in the component
of $x$ and is killed by $mp^a$.

The coordinates of $t$ are algebraic over $k$.  For sufficiently large
$r$, the point $t'=F^r(t)$ of $P^{(p^r)}$ is defined over $k^{\mathrm{sep}}$.
The finite Galois orbit of $t'$ generates a finite group, since the
group is commutative and every generator has the same finite-order
bound.  This Galois-stable set of geometric points defines a finite
\'etale $k$-subgroup $D\subset P^{(p^r)}$.  The inverse image
$H=(F^r)^{-1}(D)$ is a finite $k$-subgroup of $P$ and contains $t$
after extension to $\bar k$.  By hypothesis $H\subset Q$, so
$Q_{\bar k}$ meets the component of $x$.  Being open and closed and
containing the identity component, it therefore equals $P_{\bar k}$.
Faithfully flat descent gives $Q=P$.
\end{proof}

\begin{remark}[Related finite-group criteria]
For related criteria detecting finite-type commutative groups,
\cite[Lemmas 3.4.1--3.4.3, pp.~14--16 of the arXiv version]
{CotnerZavyalov2024B} establish lifting of primary torsion to
quotients, torsion-density statements in positive characteristic,
and an isomorphism criterion that also tests all finite $p$-power
group schemes.  Those statements likewise distinguish finite
connected subgroup schemes from their geometric points.  We have
given the finite-subgroup argument above with the hypotheses and
field descent needed for the present application.
\end{remark}

\subsubsection*{Statement B.1g (Proved): normal source and regular target over any field}
Let $k$ be a field of characteristic $p>0$, and let $f:X\to Y$ be a
proper birational morphism of proper geometrically integral
$k$-schemes.  If $X$ is normal and $Y$ is regular, then
\[
 f^*:\operatorname{Pic}^{\tau}_{Y/k}
       \xrightarrow{\ \sim\ }\operatorname{Pic}^{\tau}_{X/k}.
\]
Geometric normality of $X$ or $Y$ is not required in this statement.
\begin{proof}
The two Picard torsion components are group schemes of finite type
\cite[Corollary 4.18.3, Proposition 6.12 and Remarks 6.14 and 6.19]{Kleiman2005}.
Normality of $Y$ gives $f_*\mathcal O_X=\mathcal O_Y$ universally over
$k$.  The projection formula consequently makes $f^*$ a monomorphism;
a monomorphism of algebraic groups is a closed immersion
\cite[Theorem 5.34]{MilneAG2022B}.  Write its image as $Q\subset P$.

Geometric integrality implies that $X$ has a nonempty smooth locus.
A point in this locus is defined over a finite separable extension
of $k$.  We may make this extension, since normality and regularity
are preserved and an isomorphism of the Picard schemes descends.
We then have rational points of $X$ and $Y$ compatible with $f$.
For every finite commutative $k$-group $H$, Raynaud's correspondence
and the sections give
\[
 \operatorname{Hom}(H,\operatorname{Pic}_{Z/k})
   =H^1(Z,H^D)/H^1(k,H^D),\qquad Z=X,Y,
\]
as in B.1e.  Finite-flat torsor purity B.1b makes these groups
isomorphic.  Each homomorphism from $H$ lands in $\operatorname{Pic}^
{\tau}$, because its component-group image is torsion.  In particular
every finite $k$-subgroup scheme of $P$ factors through $Q$.
Statement B.1f now gives $Q=P$.
\end{proof}


\subsubsection*{Statement B.2h (Proved): full stabilization over a perfect field}
Let $k$ be perfect of characteristic $p>0$, and let $K/k$ be finitely
generated.  The direct system of $\operatorname{Pic}^{\tau}_{X/k}$ over
all proper normal models of $K/k$ is essentially constant.
\begin{proof}
Let $k_0$ be the algebraic closure of $k$ in $K$.  It is a finite
separable extension.  Every proper normal model has its canonical
$k_0$-structure because its ring of global functions is $k_0$, and
every domination preserving the marking by $K$ is $k_0$-linear.
An isomorphism of the Picard schemes over $k_0$ gives an isomorphism
over $k$, as can be checked after extending to $\bar k$ and taking
the product over the embeddings of $k_0$.  We may therefore assume
that $k$ is algebraically closed in $K$.

By B.2g choose a normal projective model $S$ of $K\bar k/\bar k$
after which the geometric system is constant.  Its equations and
function-field marking descend to a finite extension of $k$.  Since
$k$ is perfect, enlarge this to a finite Galois extension $l/k$.
Write $S_l$ for the descended model and $\Delta=\operatorname{Gal}(l/k)$.
Take the closure of the rational diagonal with function field $Kl$
in the product of the $\Delta$-conjugates of $S_l$, and normalize it.
The resulting projective normal $l$-scheme $N_l$ has the canonical
semilinear $\Delta$-action and dominates every conjugate of $S_l$.
Its descent datum is effective by
\cite[Lemma 39.25.3, Tag 0CCJ]{StacksB}; equivalently, one can descend
a tensor product of the conjugate ample line bundles and its graded
section algebra.  We obtain a normal projective $k$-model $N$ of $K$.
Normality is geometric here because the ground fields are perfect.

For every normal proper $W\to N$, the base change $W_{\bar k}$ is
normal and dominates $S$.  Thus the map
$\operatorname{Pic}^{\tau}_{N/k}\to\operatorname{Pic}^{\tau}_{W/k}$
becomes an isomorphism over $\bar k$ and is therefore an isomorphism.
Normalized graphs give a common normal model of any two proper
models over the perfect field.  Consequently the models dominating
$N$ form a cofinal subsystem, and the isomorphisms just obtained
give essential constancy in the usual filtered sense.
\end{proof}

\subsubsection*{Statement B.2i (Refuted): filteredness over an imperfect field}
The category of marked proper geometrically normal models of a fixed
function field need not be filtered under domination.  This failure
already occurs for two smooth projective rational threefolds over an
imperfect field.  It invalidates a normalized-graph cofinality argument
in that generality; it does not itself refute a Picard-constancy
assertion formulated for a nonfiltered category.
\begin{proof}
Let $k=\mathbb F_p(a,b)$, with $a,b$ algebraically independent, and
let $K=k(x,y,z)$.  Take $X=\mathbb P^2_k\times\mathbb P^1_k$ with
its standard marking.  Take a second copy $Y$ with the marking given
on dense affine charts by
\[
 (x,y,z)\dashrightarrow (x,y,w),\qquad
 w=z\frac{x^p-a}{y^p-b}.
\]
The inverse multiplies $w$ by $(y^p-b)/(x^p-a)$, so the map is
birational.  Both models are smooth and projective over $k$.
Put $f=x^p-a$ and $g=y^p-b$.  On the open where $z,w$ are invertible,
the closure of the graph has coordinate ring
\[
 R=B[w,w^{-1}],\qquad
 B=k[x,y,t,t^{-1}]/(y^p-b-t(x^p-a)),\quad t=z/w.
 \tag{B.2.5}
\]
The polynomial in $y$ is irreducible over $k(x,t)$, since
$b+t(x^p-a)$ is not a $p$-th power: its $t$-derivative is $x^p-a\ne0$.
Thus these rings are domains.

We verify that $B$, and hence $R$, is regular over $k$.  Where
$f\ne0$, the derivative of the defining equation with respect to
$t$ is nonzero, so the hypersurface is smooth.  On $f=0$, its
quotient is
\[
 B/(f)=l[t,t^{-1}],\qquad l=k(a^{1/p},b^{1/p}),\quad [l:k]=p^2.
\]
The element $f$ is a nonzerodivisor.  At every prime above $(f)$,
the local ring modulo $f$ is regular; lifts of its regular parameters,
together with $f$, generate the maximal ideal in the dimension
required for regularity.  This proves the assertion.  In particular
the normalization of the entire graph leaves this open unchanged.

After extending to $\bar k$ and putting
$u=x-a^{1/p}$ and $v=y-b^{1/p}$, the same open becomes
\[
 \bar k[u,v,t,t^{-1},w,w^{-1}]/(v^p-tu^p).
\]
At the codimension-one prime $(u,v)$, the local ring has dimension
one and its maximal ideal needs the two generators $u,v$; the
relation has degree at least two.  It is not regular, and therefore
is not normal.  Explicitly, its normalization adjoins $s=v/u$, with
$s^p=t$, and is $\bar k[u,s,s^{-1},w,w^{-1}]$.

Suppose a proper geometrically normal model $Z$ dominated both
marked models.  Its map to $X\times_kY$ factors through the graph
and then through the normalization of the graph, by normality.
This is a proper birational morphism.  It is an isomorphism over
an open containing the generic point of the divisor $f=0$ in
(B.2.5), because the target is normal.  Consequently $Z$ has the
same geometrically nonnormal local open near that divisor,
contradicting geometric normality.  No common such refinement exists.
The accompanying symbolic certificate verifies the graph equation,
the inverse rational map, and the normalization substitutions for
$p=2,3,5$; the proof applies to every prime.
\end{proof}

\subsection{B.3: a bound without fundamental-group theory}

\subsubsection*{Statement B.3 (Proved)}
For every finitely generated extension $K/k$ of characteristic
$p>0$ there is an integer $n(K/k)$ with
\[
 \dim\operatorname{Pic}_{X/k}\leq n(K/k)
 \tag{B.3.1}
\]
for every proper geometrically normal model $X$.  We may take
$n(K/k)=b(K/k)$ from B.2.  This proof uses neither fundamental
groups nor properness of $\operatorname{Pic}^{\tau}$.
\begin{proof}
The identity tangent space of the Picard scheme is
$H^1(X,\mathcal O_X)$
\cite[Theorem 5.11 and Corollary 5.13]{Kleiman2005}.
The dimension of a local ring is at most its embedding
dimension.  Translation in a group scheme, after algebraic
closure if necessary, shows that every component of the
Picard scheme has the identity component's dimension.
Hence
$\dim\operatorname{Pic}_{X/k}\leq h^1(X,\mathcal O_X)$,
and B.2 gives (B.3.1).  In the proof of B.2 we used the
Artin--Schreier sequence, finite covers, and a geometric
irreducibility theorem for hyperplane sections.  We never
introduced a fundamental group or used its structure or
finite generation.  Properness of a Picard group scheme
enters B.1, but no assertion from B.1 is used in B.2 or in
this dimension argument.
\end{proof}

\subsubsection*{Statement B.3a (Proved correction of the rank formula)}
Let $X$ be proper and geometrically normal over an algebraically
closed field of characteristic $p>0$, and let
$A=(\operatorname{Pic}^{0}_{X/k})_{\mathrm{red}}$, of dimension
$a$.  For every prime $\ell\ne p$,
\[
 \operatorname{rank}_{\mathbb Z_\ell}T_\ell A
 =2a
 =\dim_{\mathbb Q_\ell}H^1_{\text{ét}}
       (X,\mathbb Q_\ell).
 \tag{B.3.2}
\]
The analogous assertion with $\ell=p$ is false.
\begin{proof}
Multiplication by $\ell^r$ on the abelian variety $A$ is finite
\'etale of degree $\ell^{2ra}$, giving
$A[\ell^r](k)\simeq(\mathbb Z/\ell^r)^{2a}$
\cite[Tag 03RP]{StacksB}.  The Kummer
sequence, and divisibility of $k^\times$, identify
$H^1_{\text{ét}}(X,\mu_{\ell^r})$ with
$\operatorname{Pic}(X)[\ell^r]$.  The N\'eron--Severi group is
finitely generated and has finite torsion
\cite[Remark 6.19]{Kleiman2005}; passage to the
Tate module removes that finite component group.  These
observations give (B.3.2), up to the usual Tate twist, which
does not affect rank or dimension.

For a completely explicit characteristic two example take
the smooth plane cubic
$E: y^2+y=x^3$ with the point at infinity as identity over
$\bar{\mathbb F}_2$.  Its inverse map is
$(x,y)\mapsto(x,y+1)$, so no affine point equals its inverse.
Consequently $E[2](\bar{\mathbb F}_2)=\{0\}$, and by induction
$E[2^r](\bar{\mathbb F}_2)=\{0\}$ for every $r$.
The $2$-adic Tate module is zero, whereas $2\dim E=2$.
Smoothness follows directly from the homogeneous equation
$Y^2Z+YZ^2=X^3$: its partial derivatives in characteristic
two are $X^2,Z^2,Y^2$, with no simultaneous projective zero.
This supplies a smooth proper normal counterexample to the
uncorrected formula.  The finite certificate verifies the
polynomial identities and the inverse-pair property over
$\mathbb F_2$ and $\mathbb F_4$; the displayed algebra proves
the assertion over the algebraic closure.
\end{proof}

\subsection{Lead L6: deformation theory and direct-image vanishing}

\subsubsection*{Statement L6a (Proved under the stated cohomological hypotheses)}
Let $f:X\to Y$ be a morphism of proper geometrically integral
$k$-schemes.  If pullback is an isomorphism on $H^1(\mathcal O)$
and injective on $H^2(\mathcal O)$, then pullback induces an
isomorphism of the formal Picard functors at the identity.
For the morphisms in B.1, this gives the missing infinitesimal
step; L6c below proves that its $H^2$-injectivity
assumption holds automatically in every dimension.
\begin{proof}
We work with line bundles reducing to the trivial bundle,
modulo bundles from the local Artinian base.  We induct on
the length of a local Artinian $k$-algebra $R$.  For the
induction step choose a small extension
$0\to I\to R'\to R\to0$, where $I^2=0$ and the maximal
ideal of $R'$ annihilates $I$.  The multiplicative sheaf
sequence
\[
 1\longrightarrow 1+I\mathcal O_Z
 \longrightarrow\mathcal O_{Z_{R'}}^\times
 \longrightarrow\mathcal O_{Z_R}^\times
 \longrightarrow1
\]
identifies $1+I\mathcal O_Z$ with
$\mathcal O_Z\otimes_kI$ by $1+x\mapsto x$.
Thus the obstruction to lifting an invertible sheaf lies
in $H^2(Z,\mathcal O_Z)\otimes_kI$, and, if it vanishes,
the set of lifts is a torsor under
$H^1(Z,\mathcal O_Z)\otimes_kI$.  Global units lift here
because $H^0(Z_{R'},\mathcal O)=R'$, so these are exactly
the kernel and obstruction groups for the Picard functor.

Let a formal line bundle on $X_{R'}$ be given.  The
induction hypothesis supplies its unique preimage on
$Y_R$.  The obstruction to lifting this preimage maps
to zero on $X$, hence is zero by $H^2$-injectivity.
Choose a lift on $Y_{R'}$.  Its pullback differs from
the given bundle by an element of
$H^1(X,\mathcal O_X)\otimes I$, which comes from a unique
element on $Y$ by the $H^1$ isomorphism; adjusting the
lift gives existence and uniqueness.  This proves the
inductive assertion and hence the formal isomorphism.
\end{proof}

\subsubsection*{Statement L6b (Known vanishing and its checked surface consequence)}
For a proper birational morphism between smooth varieties
over any field,
$Rf_*\mathcal O_X=\mathcal O_Y$
\cite[Corollary 3.2.7, p.~47 of the preprint]{ChatzistamatiouRulling2011}.
In the situation of B.1 with $\dim X=\dim Y=2$ and $X$
only normal, we still have
$Rf_*\mathcal O_X=\mathcal O_Y$, and the singularities
of $X$ are rational.  This assertion holds in characteristic
$p>0$ as well as in characteristic zero.
\begin{proof}
The first assertion is exactly the cited corollary,
whose hypotheses are smoothness of both varieties and
proper birationality; no projectivity assumption is
missing.  We also opened the authors' erratum: its correction
to the correspondence functor leaves this corollary and
its proof unchanged \cite{ChatzistamatiouRullingErratum}.
For the second, take a resolution
$g:Z\to X$ of the normal surface.  The surface $Z$ is
smooth, and the composite $h=f\circ g:Z\to Y$ is
proper birational between smooth surfaces.  The first
assertion gives $R^ih_*\mathcal O_Z=0$ for $i>0$.
Here resolution exists by \cite[Introduction, Theorem,
p.~151]{Lipman1978}, since our surfaces are excellent.
Because $g_*\mathcal O_Z=\mathcal O_X$, the Leray
sequence begins
\[
 0\longrightarrow R^1f_*\mathcal O_X
 \longrightarrow R^1h_*\mathcal O_Z
 \longrightarrow f_*R^1g_*\mathcal O_Z
 \longrightarrow R^2f_*\mathcal O_X.
\]
Every fibre of $f$ has dimension at most one, so
$R^2f_*\mathcal O_X=0$ by the coherent cohomological
dimension bound \cite[Tag 0A4T]{StacksB}.
The middle term also vanishes;
therefore $R^1f_*\mathcal O_X=0$ and
$f_*R^1g_*\mathcal O_Z=0$.  The coherent sheaf
$R^1g_*\mathcal O_Z$ is supported on the finite singular
locus of the normal surface $X$.  Pushforward of such a
sheaf is faithful, hence the sheaf itself is zero.
This is the defining resolution criterion for rational
surface singularities.  Higher direct images vanish
by the same fibre-dimension bound.  We have therefore
verified the vanishing needed in L6 directly, without
using an unverified characteristic-zero implication
about the word ``sandwiched''.
\end{proof}

The final proposed search condition in L6, that a
counterexample would require dimension at least three,
nonzero $R^1f_*\mathcal O_X$, and nonreduced source
Picard torsion component, is superseded by B.1: under
the actual B.1 hypotheses no counterexample exists in
any dimension.  The conditional deformation argument
is valid, but its $H^2$ hypothesis is not needed by the
purity proof of the full theorem.


\subsubsection*{Statement L6c (Proved): first direct-image vanishing in every dimension}
\begin{theorem}
\label{B:L6-R1-vanishing}
Let $f:X\to Y$ be a proper birational morphism of integral
Noetherian schemes of characteristic $p>0$.  Suppose that $X$
is normal and $Y$ is regular.  Then
\[
                    R^1f_*\mathcal O_X=0.
\]
No dimension, projectivity, perfectness, or finite-type-over-a-field
assumption is required.  In particular, for proper varieties as in
B.1 or B.1g, pullback is an isomorphism on $H^1(\mathcal O)$ and
injective on $H^2(\mathcal O)$.  Thus the cohomological hypothesis
in L6a holds automatically for these morphisms.
\end{theorem}
\begin{proof}
We first record the semilinear fact needed over an imperfect residue
field.  Let $\kappa$ be separably closed of characteristic $p$, and
let $M\ne0$ be a finite-dimensional $\kappa$-vector space with an
injective additive map $F$ satisfying $F(av)=a^pF(v)$.  Put
\[
             V_n=\operatorname{span}_{\kappa}F^n(M).
\]
These are nonzero decreasing subspaces, and
$\operatorname{span}_{\kappa}F(V_n)=V_{n+1}$.  Choose $n$ after
the dimensions stabilize and put $V=V_n=V_{n+1}$.  In a basis
of $V$, the restriction of $F$ is $v\mapsto A v^{[p]}$ with
$A\in\operatorname{GL}_r(\kappa)$ and $r=\dim V>0$:
its column span is precisely $\operatorname{span}_{\kappa}F(V)=V$.
The kernel of $A v^{[p]}-v$ is finite \emph{\'etale} of degree
$p^r$.  Indeed, multiplication by $A^{-1}$ writes its equations
as $v_i^p-\sum_j(A^{-1})_{ij}v_j=0$; their pairwise relatively
prime leading monomials show that the monomials
$v_1^{e_1}\cdots v_r^{e_r}$ with $0\le e_i<p$ form a basis of
the quotient ring.  The Jacobian matrix of the original equations
is $-I_r$, so this finite scheme is \emph{\'etale}.  Since $\kappa$
is separably closed, it has $p^r>1$ rational points.  Consequently
$M$ has a nonzero element fixed by $F$.
This argument uses stabilized spans; injectivity of a semilinear
map alone does not imply that its matrix is invertible over an
imperfect field.

Suppose now that $\mathcal M=R^1f_*\mathcal O_X$ is nonzero.
It is coherent.  Choose a generic point of an irreducible
component of its support, localize $Y$ there, and then take a
strict henselization $A$.  The ring $A$ is Noetherian and regular,
and its residue field $\kappa$ is separably closed
\cite[Tags 06DI and 06LN]{StacksB}.  The resulting base change
$T\to\operatorname{Spec}A$ is proper and birational, and $T$
is normal: normality is preserved by localization and by the
ind-\'etale base change defining strict henselization.  It is
integral as well.  For example, flat base change preserves the
schematic density of the birational isomorphism locus, whose
base change is a nonempty open in the integral regular scheme
$\operatorname{Spec}A$; normality then excludes any additional
component.

By flat base change for coherent cohomology
\cite[Lemma 30.5.2, Tag 02KE]{StacksB},
\[
 H^0(T,\mathcal O_T)=A,\qquad
 M:=H^1(T,\mathcal O_T)\ne0,
\]
and $M$ is a finite-length $A$-module.  Finite length follows
from the choice of a minimal support prime; nonvanishing follows
from faithful flatness.  The identity in degree zero also follows
directly from proper birationality and normality of the base.

We apply B.1b to the finite flat groups $\alpha_p$ and
$(\mathbf Z/p\mathbf Z)_A$.  Its input applies exactly in this
scope: \cite[Proposition 3, p.~432, and Theorem 6,
pp.~433--434]{Marrama2026} require normality for full faithfulness
and regularity of the base, with complement of codimension at
least two, for torsor purity.  They impose no perfectness or
field-base hypothesis.  Hence
\[
 H^1_{\mathrm{fppf}}(A,G)
       \xrightarrow{\ \sim\ } H^1_{\mathrm{fppf}}(T,G_T)
       \quad(G=\alpha_p,\ \mathbf Z/p\mathbf Z).
\]
Quasi-coherent cohomology agrees in the Zariski and fppf
topologies \cite[Proposition 35.9.3, Tag 03DW]{StacksB}.
The sequence
\[
 0\longrightarrow\alpha_p\longrightarrow\mathbb G_a
       \xrightarrow{F}\mathbb G_a\longrightarrow0
\]
and the equality of global functions therefore identify
$H^1(A,\alpha_p)=A/F(A)$ with the same subgroup of
$H^1(T,\alpha_p)$.  The displayed torsor isomorphism forces
\[
                         \ker(F:M\to M)=0.
\]
Similarly, the Artin--Schreier sequence gives
\[
                      \ker(F-1:M\to M)=0,
\]
because every finite \emph{\'etale} torsor over the strictly
henselian local ring $A$ is trivial, so $H^1(A,\mathbf Z/p\mathbf Z)=0$.

Let $\mathfrak m$ be the maximal ideal of $A$ and choose $N$
with $\mathfrak m^N M=0$.  For $a\in\mathfrak m$, $v\in M$,
and $p^e\ge N$, semilinearity gives
\[
                    F^e(av)=a^{p^e}F^e(v)=0.
\]
The injectivity of $F$ implies $av=0$.  Thus $\mathfrak mM=0$,
so $M$ is a nonzero finite-dimensional $\kappa$-vector space.
The semilinear fact proved above now gives a nonzero fixed
point of $F$, contradicting $\ker(F-1)=0$.  This proves the
vanishing of $R^1f_*\mathcal O_X$.

Finally, the low-degree Leray sequence, together with
$f_*\mathcal O_X=\mathcal O_Y$, gives the asserted
$H^1$ isomorphism and $H^2$ injection.  L6a then supplies
the corresponding formal Picard isomorphism.
\end{proof}

\begin{remark}[Scope of the vanishing]
The proposed $R^1f_*\mathcal O_X\ne0$ counterexample route in
L6 is therefore impossible under its normal-source and
regular-target hypotheses in every dimension.  We have proved
only the first higher direct-image vanishing here; the theorem
does not assert that all $R^if_*\mathcal O_X$ vanish for a
normal source.  This vanishing alone does not control the
Picard images at singular targets; Statements B.2t--u below
supply that additional argument over arbitrary fields.
\end{remark}


% Checked support lemma for Part B; not included in paper_C.tex.
% The image bound needed for imperfect-field stabilization is not claimed.
\subsubsection*{Statement B.2j (Proved): a degree bound for alteration kernels}
\begin{lemma}[A degree bound for the alteration kernel]\label{B:B2-alteration-kernel}
Let $k$ be a field and let $f:Y\to X$ be a proper dominant
generically finite morphism of proper integral $k$-schemes.  Suppose
that $X$ and $Y$ are normal and put $d=[k(Y):k(X)]$.  Then the kernel
of the homomorphism of fppf Picard sheaves
\[
 f^*:\operatorname{Pic}_{X/k}\longrightarrow
       \operatorname{Pic}_{Y/k}
\]
is annihilated by $d$.  The assertion includes its infinitesimal
scheme structure, whenever the Picard functors are representable.
\end{lemma}
\begin{proof}
Write $Y\xrightarrow{h}Z\xrightarrow{\pi}X$ for the Stein
factorization.  The scheme $Z$ is normal: on an affine open of $X$,
its ring is the ring of regular functions on the inverse image in
the normal integral scheme $Y$, hence is integrally closed in its
fraction field.  Thus $Z$ is the finite normalization of $X$ in
$k(Y)$, and $h_*\mathcal O_Y=\mathcal O_Z$.
This last identity remains true after every $k$-base change by flat
base change for quasi-coherent cohomology
\cite[Lemma 30.5.2, Tag 02KH]{StacksB}.
The projection formula therefore makes $h^*$ a monomorphism of
relative Picard sheaves.

The finite torsion-free $\mathcal O_X$-module $\pi_*\mathcal O_Z$
has rank $d$.  At every point of codimension one the local ring of
$X$ is a discrete valuation ring, so this module is free there.
Consequently there is an open $j:V\hookrightarrow X$ whose complement
has codimension at least two and over which $\pi$ is finite locally
free of rank $d$.

We record explicitly why restriction to $V$ detects line bundles
universally.  Normality gives $\mathcal O_X=j_*\mathcal O_V$ by
the depth-two extension theorem
\cite[Lemma 31.13.12, Tag 0EBJ]{StacksB}.
Flat base change along $X_T\to X$, for any $k$-scheme $T$, gives
\[
 \mathcal O_{X_T}=(j_T)_*\mathcal O_{V_T}.
\]
For any invertible sheaf $\mathcal L$ on $X_T$, the projection
formula gives $\mathcal L=(j_T)_*j_T^*\mathcal L$.
An isomorphism of two invertible sheaves over $V_T$, and its inverse,
therefore extend uniquely to mutually inverse homomorphisms on $X_T$.
In particular $\operatorname{Pic}(X_T)\to\operatorname{Pic}(V_T)$
is injective.  No normality or reducedness of $T$ is being assumed.

Let a section of the kernel be represented, fppf locally on $T$,
by an invertible sheaf $\mathcal L$ on $X_T$ whose pullback to
$Y_T$ is trivial.  The monomorphism for $h$ shows that
$\pi_T^*\mathcal L$ is trivial on $Z_T$.  Over $V_T$, the norm
for the finite locally free map has the identity
\[
 \operatorname{Nm}_{\pi_T}(\pi_T^*\mathcal L)
       \simeq\mathcal L^{\otimes d};
\]
the norm and this identity commute with arbitrary base change
\cite[Lemmas 31.18.2 and 31.18.6, Tags 0BCY and 0BD2]
{StacksB}.  Hence $\mathcal L^{\otimes d}$ is trivial on
$V_T$, and the preceding extension argument makes it trivial on
$X_T$.  The result descends as an assertion about fppf sheaves.
\end{proof}

\subsubsection*{Statement B.2k (Proved): a uniform kernel along a fixed alteration}
\begin{corollary}[Uniform kernel along a fixed alteration]
\label{B:B2-fixed-alteration-kernel}
Assume characteristic $p>0$.  Let $X_0$ be a proper geometrically
normal geometrically integral $k$-scheme and let $Y\to X_0$ be a
proper generically finite map of degree $d$ with $Y$ regular and
geometrically integral.  For every proper geometrically normal
modification $X\to X_0$, let $Y_X$ be the normalization of the
main component of $Y\times_{X_0}X$.  Then
\[
 \operatorname{Pic}^{\tau}_{Y_X/k}
       =\operatorname{Pic}^{\tau}_{Y/k}
\]
canonically by Statement B.1g, and the kernel of the resulting map
\[
 h_X:\operatorname{Pic}^{\tau}_{X/k}\longrightarrow
             \operatorname{Pic}^{\tau}_{Y/k}
\]
is killed by $d$.  If, after a fixed algebraically closed extension,
the Picard groups of these modifications embed in a fixed
$g$-dimensional Jacobian as in Statement B.2d, then
\[
             |\ker h_X|\leq d^{2g}
\]
for every $X$.
\end{corollary}
\begin{proof}
The map $Y_X\to Y$ is proper birational with normal source and
regular target, and all schemes have geometrically integral
function fields.  Statement B.1g applies without requiring $Y_X$
to be geometrically normal.  The map $Y_X\to X$ has degree $d$,
so the lemma kills its Picard kernel by $d$.  Under the stipulated
closed immersion into the fixed Jacobian $J$, that kernel is a
closed subgroup of $J[d]$, whose rank is $d^{2g}$.
\end{proof}

\begin{remark}[The separate image problem]
These bounds alone do not settle stabilization over an imperfect
field.  They do not show that the images $h_X$ have uniformly
bounded finite quotients over the image of one initial model.
Choosing a model with maximal dimension of its abelian Picard part
does not by itself supply this bound: a fixed abelian quotient can
contain finite subgroup schemes of unbounded order.  A quotient
model furnished after a purely inseparable extension may fail
geometric normality over $k$, so the filtered normalized-graph
argument cannot be silently imported into the permitted category.
Statements B.2t--u solve this image problem by successive
restrictions at bad divisorial valuations.
\end{remark}


\subsubsection*{Statement B.2l (Proved): stabilization after finite purely inseparable constants}
\begin{proposition}
\label{B:B2-finite-constants}
Let $k$ have characteristic $p>0$ and let $K/k$ be a finitely
generated geometrically integral field extension.  There exist a
finite purely inseparable extension $k'/k$ and a geometrically normal
projective model $S$ of $Kk'/k'$ such that, for every proper
geometrically normal model $W$ with a birational morphism $W\to S$,
the pullback
\[
 \operatorname{Pic}^{\tau}_{S/k'}\longrightarrow
 \operatorname{Pic}^{\tau}_{W/k'}
\]
is an isomorphism.
\end{proposition}
\begin{proof}
Let $k^{\mathrm{perf}}$ be the perfect closure of $k$ in an algebraic
closure.  Geometric integrality ensures that
$K\otimes_k k^{\mathrm{perf}}$ is a field.  By Statement B.2h,
we may choose a projective normal model $S_{\infty}$ of this field
over $k^{\mathrm{perf}}$ above which all proper normal birational
models have the same $\operatorname{Pic}^{\tau}$.

The scheme $S_{\infty}$, a projective embedding, and the marking
of its function field by $Kk^{\mathrm{perf}}$ use only finitely
many coefficients.  Thus they descend to a finite purely inseparable
subextension $k'/k$ of $k^{\mathrm{perf}}/k$, giving a projective
model $S$ of $Kk'/k'$ whose base change is $S_{\infty}$.
To make the function-field statement explicit, we descend mutually
inverse birational maps between $S_{\infty}$ and a fixed projective
model of $K$ after base change; their identities on dense opens
also descend after enlarging $k'$ finitely
\cite[Lemma 32.10.1, Tag 01ZM]{StacksB}.
Since $S_{\infty}$ is normal over a perfect field, it is
geometrically normal.  The same follows for $S$ by the field-extension
criterion for geometric normality
\cite[Lemmas 33.10.3--4, Tags 038O and 038P]{StacksB}.

Now let $W\to S$ have the asserted properties.  Its base change
to $k^{\mathrm{perf}}$ is a proper birational morphism of normal
models.  The defining property of $S_{\infty}$ gives the required
Picard isomorphism after this faithfully flat field extension.
Formation of the relative Picard scheme and its $\tau$-part commutes
with field extension, so the isomorphism descends to $k'$.
\end{proof}

\begin{remark}[The precise descent left over]
The proposition supplies the weak upper-cone assertion after one
finite extension of constants.  It neither produces the model $S$
over the original field $k$ nor makes its upper cone cofinal among
the geometrically normal models over $k'$.  In particular it does
not erase the obstruction of Statement B.2i.  The additional
original-field realization is supplied by B.2t--u; its proof
uses semilinear alteration constants and valuation restrictions,
rather than only a finite field of definition for this model.
\end{remark}


\subsubsection*{Statement B.2m (Proved): the kernel bound with arbitrary alteration constants}
\begin{proposition}
\label{B:B2-alteration-constants}
Let $X_0$ be a proper geometrically normal geometrically integral
$k$-scheme of characteristic $p>0$.  We may choose a finite field
extension $l/k$, a smooth projective geometrically integral
$l$-scheme $Y$, and a proper generically finite $k$-morphism
$Y\to X_0$.  Put $d=[k(Y):k(X_0)]$ and
\[
 P=\operatorname{Pic}^{\tau}_{Y/l},\qquad
 \mathcal R=\operatorname{Res}_{l/k}(P),
 \quad\mathcal R(T)=P(T\times_k l).
\]
For each proper geometrically normal modification $X\to X_0$,
there is a compatible homomorphism of fppf sheaves
\[
 h_X:\operatorname{Pic}^{\tau}_{X/k}\longrightarrow\mathcal R
\]
whose kernel is killed by $d$.  There is an integer $g$ depending
only on $k(X_0)/k$ such that all these kernels are finite of rank
at most $d^{2g}$.  Thus a fixed comparison sheaf and a uniform
kernel bound are available without assuming that an alteration
is geometrically integral over the original field.
\end{proposition}
\begin{proof}
For existence, apply the alteration theorem over an algebraic
closure of $k$, and descend the smooth projective geometrically
integral source and its proper generically finite morphism to a
finite extension $l/k$, enlarging $l$ if necessary.  Smoothness,
geometric integrality and projectivity descend, and the finite
generic degree can be descended on an open where the map is finite
locally free.  This uses the projective alteration supplied by
\cite[Corollary 5.15, with the projectivity assertion in
Theorem 5.9 and Proposition 5.11]{deJong1997} and finite-presentation
descent \cite[Tag 01ZM]{StacksB}.  No separability assertion about
$l/k$ is required.

Let $Y_X$ be the normalization of the main component of
$Y\times_{X_0}X$.  The projection $Y_X\to Y$ is a proper
birational morphism of $l$-schemes.  Its source is normal and
geometrically integral over $l$: its function field is $k(Y)$,
which is a regular extension of $l$.  Statement B.1g over $l$
therefore identifies
\[
 \operatorname{Pic}^{\tau}_{Y_X/l}=P.
\]
Pullback from $X_l$ to $Y_X$, followed by this identification,
defines $h_X$.  Pullback preserves the $\tau$-part, so its target
is the indicated Weil restriction.  Formation of Weil restriction
commutes with base change, and $\mathcal R$ is an algebraic space
\cite[Lemma 97.11.2 and Proposition 97.11.5, Tag 05Y8]{StacksB}.
It is separated: restriction of scalars along a finite locally
free map preserves a closed immersion, as is checked by the
finitely many coefficient equations of its ideal, and we apply
this to the diagonal of $P$.  In particular the kernels under
discussion are closed subgroup schemes of
$\operatorname{Pic}^{\tau}_{X/k}$.

We verify that passing from relative Picard classes over $k$ to
relative classes over $l$ has not enlarged the kernel.  Work fppf
locally on a $k$-scheme $T$ so that a section of
$\operatorname{Pic}^{\tau}_{X/k}$ is represented by a line bundle
$\mathcal L$ on $X_T$.  Write $T_l=T\times_k l$.
If $h_X(\mathcal L)=0$, then the pullback of $\mathcal L$ to
$(Y_X)_T$ has zero class in the relative Picard sheaf for
$(Y_X)_T\to T_l$.  Since
$\mathcal O_{T_l}=q_*\mathcal O_{(Y_X)_T}$ universally, the
low-degree Leray sequence for $\mathbf G_m$ says that this line
bundle is pulled back from a line bundle $\mathcal M$ on $T_l$.

Such an $\mathcal M$ becomes trivial after an fppf cover of $T$,
even if $l/k$ is inseparable and $T$ is nonreduced.  Indeed let
$\pi:T_l\to T$.  The functor of generators of $\mathcal M$ as
an $\mathcal O_{T_l}$-module is an open subscheme of the vector
bundle associated with the finite locally free
$\mathcal O_T$-module $\pi_*\mathcal M$.
The open condition is that the determinant of
$\mathcal O_{T_l}\to\mathcal M$, $1\mapsto m$, is invertible.
It is smooth over $T$ and has a point over every geometric
fibre, since an invertible module over a finite product of Artin
local rings is free of rank one.  Hence it is a smooth surjective,
and therefore fppf, cover of $T$ trivializing $\mathcal M$.
Consequently $\ker h_X$ is exactly the kernel of the ordinary
$k$-relative Picard pullback along $Y_X\to X$.

The latter proper dominant map has generic degree $d$ and normal
source and target.  Statement B.2j therefore kills its entire
fppf Picard kernel by $d$.  For the rank bound, extend to an
algebraic closure and then to the fixed generic-curve field used
in B.2d.  If $X_0$ is not projective, take the normalization of
the graph between $X$ and a fixed projective normal model before
restricting to that curve.  Pullback from $X$ to the graph is a
closed immersion on $\operatorname{Pic}^{\tau}$, and B.2d embeds
the latter into one fixed $g$-dimensional Jacobian $J$.
The kernel is therefore a closed subgroup of $J[d]$ and has
rank at most $d^{2g}$.  Finiteness and this rank bound descend
through the field extensions.
\end{proof}

\begin{remark}
We have not identified the $\tau$-part of the Picard scheme of
the possibly geometrically nonreduced $k$-scheme $Y$ with a Weil
restriction.  The kernel calculation only uses the relative
Picard scheme over $l$, the explicit fppf trivialization of base
line bundles, and the norm argument of B.2j.  The separate
uniform bound for the relevant image quotients is proved in
B.2t--u by a valuation restriction argument.
\end{remark}


\subsubsection*{Statement B.2n (Proved): the rational relative surface step}
\begin{proposition}
\label{B:B2-surface-rational-step}
Let $f:W\to X$ be a proper birational morphism of proper
geometrically normal geometrically integral surfaces over a field
$k$.  If $R^1f_*\mathcal O_W=0$, then
\[
 f^*:\operatorname{Pic}^{\tau}_{X/k}
       \xrightarrow{\ \sim\ }\operatorname{Pic}^{\tau}_{W/k}.
\]
\end{proposition}
\begin{proof}
We may extend $k$ to an algebraic closure: normality is preserved,
and the higher-direct-image hypothesis is preserved by flat base
change.  The pullback of Picard schemes is a closed immersion
because $f_*\mathcal O_W=\mathcal O_X$ universally.  Leray gives
an isomorphism on $H^1(\mathcal O)$ and an injection on
$H^2(\mathcal O)$.  By L6a the formal Picard functors at the
identity are isomorphic.

We next prove surjectivity on geometric points.  Let $\mathcal L$
represent a point of $\operatorname{Pic}^{\tau}_{W/k}(k)$.
It is numerically trivial.  For a closed point $x\in X$, set
$A=\mathcal O_{X,x}$ and $W_A=W\times_X\operatorname{Spec}A$.
This proper map has fibres of dimension at most one and
$H^1(W_A,\mathcal O_{W_A})=0$.  The degree of $\mathcal L$
on each integral exceptional curve is zero.  The numerical
triviality criterion of \cite[Theorem 12.1(i), pp.~220--221]
{Lipman1969B} gives
$\mathcal L|_{W_A}\simeq\mathcal O_{W_A}$.
This trivialization spreads to a neighbourhood of $x$ by finite
presentation.  Equivalently, $f_*\mathcal L$ is locally free of
rank one near $x$ and its adjunction pullback is $\mathcal L$.
Thus these local descents glue, using $f_*\mathcal O_W=\mathcal O_X$,
to an invertible sheaf $\mathcal N=f_*\mathcal L$ on $X$.
For each integral curve on $X$, its strict transform on $W$
shows that $\mathcal N$ has degree zero.  Hence $\mathcal N$
is numerically trivial and belongs to $\operatorname{Pic}^{\tau}_{X/k}$
\cite[Theorem 6.3 and Remark 6.19]{Kleiman2005}.

Finally, a closed immersion of finite-type group schemes that
is an isomorphism on formal completions at the identity has
zero defining ideal in the identity local ring: completion is
faithfully flat, or one may use Krull intersection.  It is
therefore an isomorphism on a neighbourhood of the identity.
Translation by the points just lifted treats every geometric
closed point.  The coherent defining ideal is zero everywhere,
so the immersion is an isomorphism.  Faithfully flat descent
returns the conclusion over the original field.
\end{proof}

\subsubsection*{Statement B.2o (Proved): arbitrary-field surface stabilization}
\begin{theorem}
\label{B:B2-surface-stabilization}
Let $K/k$ be geometrically integral of transcendence degree two,
and suppose it has a proper geometrically normal model $X_0$.
There exists a proper geometrically normal modification $X_1\to X_0$
such that every proper geometrically normal modification $W\to X_1$
induces an isomorphism
\[
 \operatorname{Pic}^{\tau}_{X_1/k}
       \xrightarrow{\ \sim\ }\operatorname{Pic}^{\tau}_{W/k}.
\]
The field $k$ may be imperfect.  This is the weak upper-cone
formulation and does not assume filteredness of the allowed model
category.
\end{theorem}
\begin{proof}
Choose a regular proper resolution $\pi:Y\to X_0$, which exists
for excellent surfaces \cite[Introduction, Theorem, p.~151]
{Lipman1978}, and put
\[
                 \mathcal F=R^1\pi_*\mathcal O_Y.
\]
This coherent sheaf is supported on finitely many closed points,
so it has finite total length.

We need the regular-target vanishing over the actual field, not
only over a perfect extension.  A regular two-dimensional local
ring is rational by taking its identity resolution.  Thus
\cite[Definition 1.1 and Proposition 1.2(2), pp.~199--200]
{Lipman1969B} imply $R^1a_*\mathcal O_Z=0$ for every proper
birational normal surface $a:Z\to Y$ with $Y$ regular.
Together with $a_*\mathcal O_Z=\mathcal O_Y$, this makes
the sheaf $\mathcal F$ independent of a further regular resolution.

For any proper geometrically normal modification $b:X\to X_0$,
resolve the normalization of the graph of $X\dashrightarrow Y$.
We obtain a regular surface $Z$ with proper birational maps to
both $X$ and $Y$.  Leray and the preceding vanishing give a
canonical injection
\[
        R^1b_*\mathcal O_X\hookrightarrow\mathcal F.
        \tag{B.2.6}
\]
These injections are compatible with domination: two choices of
$Z$ are compared by a common regular resolution, on which all
pullback identifications agree.

Choose an allowed $b_1:X_1\to X_0$ for which the total length
of $R^1(b_1)_*\mathcal O_{X_1}$ is maximal.  Such a maximum
exists among the integers between zero and the length of
$\mathcal F$; the identity model is an allowed choice.
For $g:W\to X_1$ as in the statement, (B.2.6) and maximality
make the natural map
\[
 R^1(b_1)_*\mathcal O_{X_1}
       \longrightarrow R^1(b_1g)_*\mathcal O_W
\]
an isomorphism.  Every fibre of $b_1$ has dimension at most one,
so $R^2(b_1)_*\mathcal O_{X_1}=0$
\cite[Tag 0A4T]{StacksB}.  The low-degree Leray sequence now gives
\[
                   (b_1)_*R^1g_*\mathcal O_W=0.
\]
The coherent sheaf $R^1g_*\mathcal O_W$ is supported on finitely
many closed points of $X_1$.  Pushforward is faithful on sheaves
with such support, so $R^1g_*\mathcal O_W=0$.
Statement B.2n applies and proves the claimed isomorphism.
\end{proof}

\begin{remark}[Low-dimensional scope]
For transcendence degree zero the model is fixed, and for
transcendence degree one every proper normal model is the unique
regular proper curve with that function field.  Accordingly the
weak upper-cone stabilization is now proved in dimensions at most
two over arbitrary fields whenever the stipulated models exist.
If $k$ is not algebraically closed in $K$, we first replace it by
its finite constant field $k_0\subset K$.  Existence of a
geometrically normal model makes $k_0/k$ separable.  All marked
models and their maps are canonically $k_0$-schemes; applying the
result over $k_0$ and checking the finitely many geometric
components after an algebraic closure gives the same assertion
over $k$.
No common geometrically normal refinement was used: the common
regular surfaces in the proof were auxiliary objects only.
\end{remark}


\subsubsection*{Statement B.2p (Refuted auxiliary lemma): a graph with an elliptic target}
\begin{proposition}
\label{B:B2-abelian-graph-counterexample}
There is a geometrically normal projective integral surface $M$
over an imperfect field $k$, an elliptic curve $E/k$, and a rational
map $M\dashrightarrow E$ for which no proper geometrically normal
modification of $M$ resolves the map.  In particular, the normalization
of a graph with an abelian target need not be geometrically normal.
\end{proposition}
\begin{proof}
We work in characteristic three and put $k=\mathbf F_3(a)$, where
$a$ is transcendental.  In $\mathbf P^2_k$ let
\[
 E:\quad F:=Y^2Z-X^3+XZ^2=0,\qquad Q:=Y^2+Z^2.
\]
The gradient of $F$ is
$(Z^2,-YZ,Y^2-XZ)$.  Its simultaneous vanishing on $F=0$
would first give $Z=0$, then $X=0$, then $Y=0$, which is
impossible in projective space.  Thus $E$ is a smooth plane cubic.
With origin $[0:1:0]$ it is an elliptic curve.  The divisor
$E\cap\{Q=0\}$ consists of six distinct geometric points.
Indeed none lies on $Z=0$, and in the chart $Z=1$ its equations
are
\[
                    y^2+1=0,\qquad x^3-x+1=0.
\]
Both polynomials are separable, and their degrees are two and
three.  The gradients of $F$ and $Q$ are independent at all
these points: the first has nonzero $X$-component $Z^2$,
whereas the second has zero $X$-component and nonzero
$Y$-component $-Y$.

Let $C\subset\mathbf P^3_k$, with coordinates $[X:Y:Z:W]$,
be the projective cone $F=0$.  Define $M\subset\mathbf P^4_k$
by
\[
          F=0,\qquad V^3=aW^3+WQ.
          \tag{B.2.7}
\]
Projection forgetting $V$ is a finite flat map $M\to C$ of
degree three: on each standard affine chart of $C$ the equation
for the appropriately scaled $V$ is monic of degree three.
The projection centre does not meet $M$.

We verify geometric normality directly.  Over an algebraic closure
$\bar k$, choose $b^3=a$ and write $V'=V-bW$.  The equations
become
\[
                         F=0,\qquad (V')^3=WQ.
\]
The coordinate ring of the affine cone over $E_{\bar k}$ is a
domain, and $WQ$ is not a cube in its rational function field:
its $W$-adic valuation is one.  Hence the latter equation is
irreducible over that function field.  Since it is monic, the
homogeneous coordinate ring of $M_{\bar k}$ is a domain.
The two displayed equations form a regular sequence, so
$M_{\bar k}$ is Cohen--Macaulay and in particular satisfies $S_2$.

At a point with $(X,Y,Z)\ne(0,0,0)$, the gradient of $F$ is
nonzero.  If $Q\ne0$, the $W$-component of the gradient of
$(V')^3-WQ$ is nonzero, so the two gradients are independent.
If $Q=0$, we have $Z\ne0$ and $Y\ne0$.  A proportionality
between the gradients would have proportionality factor zero,
by their $X$-components; it therefore forces $W=0$ and then
$V'=0$.  These are precisely the six points just described.
The remaining possible singular point is
$[0:0:0:1:0]$.  Thus the geometric singular locus is finite.
As $M_{\bar k}$ has dimension two, it is regular in codimension
one.  Serre's criterion now proves that $M$ is geometrically
normal, as asserted.

For completeness we describe the graph without invoking an
unspecified resolution.  The incidence variety
\[
 Y_0=\{(c,e)\in C\times E:
              (X,Y,Z)\text{ is proportional to the coordinates of }e\}
\]
is a ruled surface over $E$, isomorphic to
$\mathbf P_E(\mathcal O_E\oplus\mathcal O_E(1))$
with the quotient convention for projective bundles.  Its map to
$C$ is an isomorphism away from the cone vertex and contracts
one section $E_0$ to that vertex.  The displayed incidence
condition consists of the three $2\times2$ minors, so $Y_0$
is a closed subscheme of $C\times E$.
Put
\[
                         \Gamma=M\times_CY_0.
\]
It is finite flat over $Y_0$ and integral, by the same monic
generic-field argument used above.  It is a closed subscheme
of $M\times E$, and is proper birational over $M$.

Near the exceptional section choose a local trivialization of
$\mathcal O_E(1)$ and a transverse parameter $t$, so that the
cone coordinates are $t$ times the coordinates on $E$ and
$W=1$.  The local equation of $\Gamma$ is
\[
                         v^3=a+t^2q,
                         \tag{B.2.8}
\]
where $q$ represents the section $Q$ in this trivialization.
The equation $t=0$ cuts out the Cartier divisor
$D=E_l$, where $l=k(a^{1/3})$.  This is a regular scheme:
$E_l$ is smooth over the field $l$, including at its closed
points whose residue fields may be inseparable over $k$.
Since $t$ is a nonzerodivisor, regularity of the quotient by
$t$ implies regularity of $\Gamma$ along $D$.  Away from $D$,
the map $\Gamma\to M$ is an isomorphism.  Therefore $\Gamma$
is normal everywhere.  We do not claim that it is regular
at the six other singular points.

Because $\Gamma$ is integral and contains the graph over the
dense open on which the projection $M\dashrightarrow E$ is
defined, its closed incidence embedding identifies it with
the graph closure.  It is already normal, so it is also the
normalized graph.

After extension to $\bar k$, the local equation at the generic
point of the reduced exceptional divisor is
\[
            s^3=t^2q,\qquad q\in\bar k(E)^\times,
            \quad s=v-b.
\]
Its local ring at $(s,t)$ is a one-dimensional local domain
of embedding dimension two, hence is not a discrete valuation
ring and is not normal.  One can also see the defect explicitly:
$h=s^2/t$ is integral because $h^3=tq^2$, and
$t=h^3/q^2$, $s=h^2/q$.  The cusp ring has value semigroup
generated by two and three, whereas $h$ has value one, so
$h$ does not belong to it.

Finally, suppose a proper geometrically normal modification
$N\to M$ resolved the rational map to $E$.  The resulting
map $N\to M\times E$ factors through the graph $\Gamma$;
it is proper and birational.  At the generic point of the
$k$-divisor $D$, the scheme $\Gamma$ is regular of dimension
one.  A proper birational morphism between normal integral
schemes is an isomorphism over a neighbourhood of such a
codimension-one point.  Thus an open neighbourhood in $N$
would be isomorphic to an open neighbourhood of that point
in $\Gamma$.  After passage to $\bar k$ this neighbourhood
contains the nonnormal codimension-one ring just computed,
contrary to geometric normality of $N$.

The accompanying exact certificate checks the homogeneous
derivatives, the affine singular-locus ideal, the six distinct
points over $\mathbf F_{3^6}$, and the integral-element
identities.  The regularity and proper-birational arguments
above supply the geometric conclusions.
\end{proof}

\begin{remark}[What this counterexample resolves]
The ordinary normal graph always exists, and over a perfect
field its normality is geometric normality.  The extra hypothesis
that the target is an abelian variety does not provide this
conclusion over an imperfect field.  Thus the proposed abelian
graph construction cannot close B.2 in its original-field scope.
This proposition is not a counterexample to weak upper-cone
Picard stabilization; indeed B.2o proves that weak assertion for
the surface in (B.2.7).
\end{remark}


\subsubsection*{Statement B.2q (Refuted): a constant cofinal Picard diagram over an imperfect field}
\begin{theorem}
\label{B:B2-cofinal-counterexample}
There is an imperfect field $k$ of characteristic $3$, a regular
function field $K/k$ of transcendence degree three, and two proper
geometrically normal models $G_1,G_2$ of $K/k$ with the following
property.  For every proper geometrically normal modification
$W\to G_i$, the geometric reduced identity component of
$\operatorname{Pic}_{W/k}$ is isomorphic to an elliptic curve
$E_i$.  The curves $E_1$ and $E_2$ are not geometrically isogenous.

Consequently there is no domination-cofinal collection of the
proper geometrically normal models of $K/k$ on which all the
$\operatorname{Pic}^{\tau}$ are isomorphic to one fixed group
scheme.  In particular a constant cofinal Picard diagram, with
comparison maps given by birational pullback, does not exist.
This does not refute the weak upper-cone formulation.
\end{theorem}

\begin{proof}
We give the construction and the Picard comparison separately.

\par\smallskip\noindent\textbf{The two elliptic factors and the generic section.}\quad
Over $\mathbf F_3$ put
\[
 E_1:\ y^2=x^3-x,\qquad
 E_2:\ y^2=x^3+x^2+1,
 \qquad A=E_1\times E_2.
\]
Both plane cubics are smooth.  Write
$L=\mathcal O_{E_1}(1)\boxtimes\mathcal O_{E_2}(1)$.
This gives a projectively normal Segre embedding of $A$.
For completeness, each cubic section ring is its normal
homogeneous hypersurface coordinate ring: the affine cone is
regular off its vertex and satisfies $S_2$.  Their Segre product
is normal, since it is the invariant subring of their normal
tensor product for the action having weights $1$ and $-1$.
It is generated in degree one.  These arguments remain valid
after every field extension.

On each plane cubic, the six quadratic monomials
$X^2,XY,XZ,Y^2,YZ,Z^2$ form a basis of
$H^0(E_i,\mathcal O_{E_i}(2))$.  Let
$q_1,\ldots,q_{36}$ be their pairwise external products,
a basis of $H^0(A,L^2)$, and set
\[
 k=\mathbf F_3(a,c_1,\ldots,c_{36}),\qquad
 Q=\sum_{r=1}^{36}c_rq_r,
 \qquad l=k(a^{1/3}),
\]
where all the displayed variables are algebraically independent.
The divisor of $Q$ on $A_k$ is geometrically smooth.  One can
check the generic-section assertion directly: very ampleness of
$L^2$ makes vanishing of the value and the two first derivatives
at a geometric point three independent linear conditions on a
section.  The incidence of singular sections thus has dimension
at most $34$ in the product with the $35$-dimensional projective
linear system.  Its image is closed because $A$ is proper, and
does not contain the generic section.  The integer $36$ follows
from $h^0(E_i,\mathcal O_{E_i}(2))=6$ and the product formula.

\par\smallskip\noindent\textbf{The cone and its purely inseparable triple cover.}\quad
Let $C$ be the projective cone over $(A_k,L)$, with cone coordinate
$W$ and vertex $o$.  Its homogeneous coordinate ring is
$S_A[W]$, where $S_A=\bigoplus_{n\ge0}H^0(A_k,L^n)$.
Define
\[
 M=\operatorname{Proj}
 \frac{S_A[W,V]}{(V^3-aW^3-WQ)}.
 \tag{B.2.9}
\]
Forgetting $V$ gives a finite flat map $M\to C$ of degree three,
with algebra
$\mathcal O_C\oplus\mathcal O_C(-1)\oplus\mathcal O_C(-2)$.
The unique point of $M$ over $o$ has residue field $l$; denote it
by $P$.

The threefold $M$ is geometrically integral and geometrically
normal.  Indeed, after adjoining $a^{1/3}$, the change
$V'=V-a^{1/3}W$ gives $V'^3=WQ$.  In the function field of the
cone the right side is not a cube, as its order in the variable
$W$ is one.  This proves geometric integrality.
The cone is geometrically normal and hence geometrically $S_2$;
the finite flat triple cover has zero-dimensional
Cohen--Macaulay fibres, so $M$ is geometrically $S_2$.
Away from $P$, the cone is smooth.  There the only singularities
of the displayed equation after an algebraic closure occur at
$W=V'=Q=0$.  Since $Q=0$ is smooth on the surface $A$, this is
a curve, of codimension two in $M$; the point $P$ has codimension
three.  Thus the geometric $R_1$ condition also holds, and
Serre's criterion proves geometric normality.  The commutative
algebra inputs for $S_2$ under field extension and flat maps are
\cite[Tags 045M, 0339 and 031S]{StacksB}.

\par\smallskip\noindent\textbf{The normal graph and its bad divisor.}\quad
Let $p:R\to A_k$ be the ruled resolution of the cone, with
exceptional section $D_0$ of normal bundle $L^{-1}$, and let
$\rho:R\to C$ be the contraction.  Put
$H=\rho^*\mathcal O_C(1)$, so that $H=D_0+p^*L$.
Form the finite flat pullback
\[
                    \Gamma=M\times_C R.
\]
The incidence description of the ruled resolution embeds
$R$ as a closed subscheme of $C\times A_k$; hence $\Gamma$
is a closed subscheme of $M\times A_k$.
It is integral and isomorphic to $M$ away from $P$.
Its exceptional divisor is
\[
                       D=A_l.
\]
Locally along $D_0$, with parameter $t$ normal to the section
and a trivialization $q$ of $Q$, its equation is
\[
                         v^3=a+t^2q.                 \tag{B.2.10}
\]
The quotient by $t$ is the regular scheme $A_l$.
Since $t$ is a nonzerodivisor, this shows that $\Gamma$ is
regular at every point of $D$.  Away from $D$ it is the normal
scheme $M\setminus\{P\}$.  Thus $\Gamma$ is normal, and its
closed incidence embedding identifies it with the normal graph
of the rational map $M\dashrightarrow A_k$.

It is not geometrically normal.  At the generic point of
$D_{\bar k}$, equation (B.2.10) becomes
\[
                    v'^3=t^2q,\qquad q\ne0.
\]
This is a one-dimensional local domain of embedding dimension
two.  It is therefore not a discrete valuation ring and is not
normal.  A proper birational morphism between normal schemes is
an isomorphism over a neighbourhood of each codimension-one
point of the target.  Consequently no geometrically normal
proper model over $M$ can resolve $M\dashrightarrow A_k$:
such a resolution would factor through its closed graph
$\Gamma$ and would retain this geometrically nonnormal
codimension-one local ring.  The same obstruction persists
after every finite separable extension of $k$, because
$a^{1/3}$ remains outside such an extension.

\par\smallskip\noindent\textbf{Two geometrically normal partial graphs.}\quad
For $i=1,2$, let $G_i$ be the normalization of the image of
\[
       \Gamma\longrightarrow M\times E_i,
\]
where the second component is the appropriate projection of
$A=E_1\times E_2$.  Normality of $\Gamma$ gives a proper
birational map $r_i:\Gamma\to G_i$.  Off $D$ it is an
isomorphism.  It contracts $D=A_l$ onto the curve $(E_i)_l$;
its fibres on $D$ are the other elliptic factor over $l$.
Indeed, the reduced image curve in $G_i$ is finite over
$(E_i)_l$, while the projection $A_l\to(E_i)_l$ has direct
image of its structure sheaf equal to $\mathcal O_{(E_i)_l}$.
Its Stein factorization therefore forces this finite
intermediate curve to be $(E_i)_l$ itself.
In particular $G_i\to M$ has no
exceptional divisor.

Every codimension-one point of $G_i$ therefore has a
neighbourhood isomorphic to an open of $\Gamma\setminus D$,
which is geometrically normal.  Thus $G_i$ is geometrically
$R_1$.  It is geometrically $S_2$ because it is normal:
field extensions have Cohen--Macaulay fibres, and $S_2$
ascends by the cited flat-map criterion.  Its function field
is the regular extension $k(M)/k$, so its geometric generic
fibre is integral.  Serre's criterion now proves that both
$G_i$ are geometrically normal.  They are projective, being
finite over closed subvarieties of $M\times E_i$.

\par\smallskip\noindent\textbf{The connected Picard scheme of the graph.}\quad
For the degree-three map $\pi:\Gamma\to R$ we have
\[
 \pi_*\mathcal O_\Gamma
 =\mathcal O_R\oplus\mathcal O_R(-H)\oplus\mathcal O_R(-2H).
\]
The projective-bundle calculation gives
\[
 p_*\mathcal O_R=\mathcal O_A,\quad R^1p_*\mathcal O_R=0,
 \qquad Rp_*\mathcal O_R(-H)=0,
\]
and
\[
 p_*\mathcal O_R(-2H)=0,\qquad
 R^1p_*\mathcal O_R(-2H)=L^{-1}.
\]
The last identity also follows from relative duality, since
$\omega_{R/A}=\mathcal O_R(-2H)\otimes p^*L$.
As $H^0(A,L^{-1})=0$, we conclude that
\[
 H^1(\Gamma,\mathcal O_\Gamma)=H^1(A,\mathcal O_A),
 \qquad h^1(\Gamma,\mathcal O_\Gamma)=2.
\]
The map $\Gamma\to A$ has universal direct image of its
structure sheaf equal to $\mathcal O_A$, so pullback embeds
$A^\vee$ as a closed subgroup of $\operatorname{Pic}^0_{\Gamma/k}$.
This subgroup has dimension two, equal to the tangent dimension
of the ambient connected group scheme.  The latter is therefore
smooth of dimension two and equals $A^\vee$.

We also need a regular birational model with the same connected
Picard scheme.  The singularities of $\Gamma$ off $D$ are
\'etale locally the product of the surface
$xy=z^3$ with a smooth curve: replace $q$ locally by
$q+aW^2$ before using the equation.  Blowing up its reduced singular
locus resolves these singularities.  On the three charts of
the blowup of $(x,y,z)$ the strict transforms are respectively
$y_1=xz_1^3$, $x_1=yz_1^3$, and $z=x_1y_1$, hence smooth.
The exceptional surface-resolution fibre is a chain of two
projective lines of self-intersection $-2$.
For its reduced union $E$, the restriction of
$\mathcal O_E(-nE)$ has degree $n$ on each component for
$n\ge0$.  The normalization exact sequence at their one
intersection gives $H^1(E,\mathcal O_E(-nE))=0$.
Induction on infinitesimal neighbourhoods and formal functions
give $R^1$ of this surface resolution equal to zero; flat and
\'etale base change give the same conclusion for the local
products.  We thus obtain a proper regular model
$s:Y\to\Gamma$, unchanged near $D$, with
$R^1s_*\mathcal O_Y=0$.  Consequently
\[
                    \operatorname{Pic}^0_{Y/k}=A^\vee
                    \tag{B.2.11}
\]
by the same tangent-dimension argument.

Pullback along $r_i$ is a closed immersion of Picard schemes.
Its image in $\operatorname{Pic}^0_\Gamma=A^\vee$ contains
$E_i^\vee$, because $G_i$ maps to $E_i$.  On the other hand,
its restriction to every contracted fibre of
$A_l\to(E_i)_l$ is trivial.  On $A^\vee$ this condition is
exactly the vanishing of the other projection, after the
faithfully flat extension $l/k$.  It follows scheme-theoretically
that
\[
                    \operatorname{Pic}^0_{G_i/k}=E_i^\vee.
                    \tag{B.2.12}
\]

\par\smallskip\noindent\textbf{Every permitted domination stays on its elliptic branch.}\quad
Let $W\to G_i$ be any proper geometrically normal modification.
Take the normalization $T$ of the main graph between $W$ and
the fixed regular model $Y$.  The map $T\to Y$ has normal
source and regular target, so B.1g identifies
$\operatorname{Pic}^{\tau}_T$ with
$\operatorname{Pic}^{\tau}_Y$.  Pullback from $W$ is a closed
immersion.  We obtain compatible closed immersions
\[
       E_i^\vee\ \subseteq\ \operatorname{Pic}^0_{W/k}
              \ \subseteq\ A^\vee.                 \tag{B.2.13}
\]
Suppose the middle group has dimension two.  A closed subgroup
of the abelian surface with that dimension and connected
underlying space is the whole abelian surface.  Hence the
second immersion is an isomorphism over $k$.

To turn this equality into the needed morphism, we may make a
finite separable field extension and choose a rational point
in the common smooth birational locus of $T,Y,W$.  Their
Albanese torsors then acquire compatible origins.  For proper
geometrically integral schemes the Albanese variety is dual to
the maximal abelian subvariety of the Picard scheme
\cite[Theorem 3.2 and Remark 3.3, pp.~12--13 of the author
manuscript]{C_ACMV2026}.  The two pullback isomorphisms to
$\operatorname{Pic}^0_T=A^\vee$ therefore dualize to
isomorphisms of Albanese varieties.  The morphism
$T\to Y\to A$ factors through a morphism $W\to A$, up to
the translation fixed by the chosen origin.  This extends the
original rational map $M\dashrightarrow A$ on $W$.
It contradicts the bad-divisor obstruction, which survives
this separable extension.  Thus the middle group in (B.2.13)
has dimension one.  After an algebraic closure its reduced
identity component is exactly $E_i^\vee$.

\par\smallskip\noindent\textbf{The two branches cannot have one constant value.}\quad
We finish by checking the distinction between the elliptic
curves without an unverified Hasse-invariant calculation.
The chord-and-tangent law
\cite[Chapter I, Section 3, Theorem 3.1, pp.~25--27]{MilneEC2021B}
says that a finite point of a smooth plane cubic is killed by
$3$ precisely when its tangent is a flex.  For
$y^2=x^3+a_2x^2+a_4x+a_6$ in characteristic three and $y\ne0$,
substitution of the tangent line shows that this condition is
$m^2=a_2$, where $m$ is its slope.  On $E_1$ we have
$a_2=0$ and $m=-1/(2y)\ne0$, so there is no such finite point;
the points with $y=0$ are of order two.  On $E_2$ we have
$m=x/y$, so the flex condition gives $y^2=x^2$, hence
$x^3+1=0$.  The two solutions are $(-1,1)$ and $(-1,-1)$.
Thus
\[
 \#E_1[3](\bar k)=1,\qquad \#E_2[3](\bar k)=3.
\]
Surjectivity of multiplication by $3$ on geometric points
gives respectively $1$ and $3^r$ points killed by $3^r$.
If the two curves were isogenous, dualizing if necessary would
give an isogeny $E_2\to E_1$.  It would kill every geometric
$3^r$-torsion point of $E_2$, placing $3^r$ distinct points in
one fixed finite kernel for every $r$, which is impossible.
The two curves are not geometrically isogenous.
The accompanying exact program verifies the cubic smoothness,
tangent substitution and geometric flex equations; it keeps
the nonreduced cube equation distinct from its two reduced
finite flex points.

Finally, a domination-cofinal collection must contain a model
dominating $G_1$ and a model dominating $G_2$.  An isomorphism
between their $\operatorname{Pic}^{\tau}$ would induce an
isomorphism of their geometric reduced identity components,
contradicting the preceding calculation.  This proves the
asserted failure of the strong cofinal-diagram formulation.
\end{proof}

\begin{remark}[Scope]
We have not asserted that the entire torsion Picard group is
constant on either upper cone.  Finite infinitesimal subgroup
schemes in the other elliptic factor may still occur in
(B.2.13).  The reduced identity components already distinguish
the two cofinal requirements and suffice for the counterexample.
\end{remark}


\begin{remark}[The two assertions in the prompt]
The cohomological assertion in B.2 is true over every ground field.
The asserted equivalent formulation as a constant domination-cofinal
Picard diagram is false by the preceding theorem.  Moreover B.2i shows
that the indicated model category need not be filtered.  A weaker
replacement asks only for one model above which all permitted pullbacks
are isomorphisms, without cofinality in the entire category.  This is
proved over every ground field in B.2u below.  Thus both
cohomological boundedness and weak upper-cone stabilization
hold, while constant cofinality fails in the present example.
\end{remark}


\subsubsection*{Statement B.2r (Proved): geometric torus exclusion and bounded exponent}
\begin{proposition}\label{B:B2-no-torus-exponent}
Let $k$ be a field of characteristic $p>0$.
\begin{enumerate}
\item If $Y$ is a proper normal $k$-variety, then
$(\operatorname{Pic}^{\tau}_{Y/k})_{\bar k}$ contains no nontrivial
torus.  Geometric normality of $Y$ is unnecessary.
\item Let $H$ be a commutative $k$-group scheme of finite type such that
$H_{\bar k}$ contains no nontrivial torus.  Then $H$ has a largest
abelian subvariety $A_H$, and $H/A_H$ is annihilated by one positive
integer.  More precisely, for all sufficiently large $n$,
\[
 A_H=\operatorname{im}([p^n]:H^0\longrightarrow H^0).
\]
\item If $l/k$ is any finite field extension and $P$ is a
proper commutative $l$-group scheme of finite type, then
$\operatorname{Res}_{l/k}P$ contains no geometric torus.  Thus (ii)
applies to every closed subgroup scheme of this restriction of scalars.
\end{enumerate}
\end{proposition}
\begin{proof}
We first prove (i), including the assertion on arbitrary test schemes
needed for a statement about group schemes.  Put $\Omega=\bar k$,
$X=Y_\Omega$, and let
\[
 \nu:Z=(X_{\mathrm{red}})^{\mathrm{norm}}\longrightarrow X.
\]
Tanaka's Lemma~2.2(2), including its proof on pp.~6--8, says that the
normalization of the reduction of any field extension of a normal
variety is a universal homeomorphism
\cite[Lemma 2.2(2)]{Tanaka2015D}.
Here normalization is finite, since $X$ is of finite type over a field.
Consequently $\nu$ is a finite universal homeomorphism.  Its kernel on
structure sheaves is nilpotent.  Kelly's Corollary~20 therefore supplies
an integer $e$ and an absolute Frobenius factorization
\[
 \sigma:X\longrightarrow Z,\qquad
 \nu\sigma=F_X^e,\qquad \sigma\nu=F_Z^e
\]
\cite[Corollary 20, p.~9]{Kelly2018D}.
The morphism $\sigma$ is $F_\Omega^e$-semilinear; it need not be an
$\Omega$-morphism.

We claim that the kernel of
$\nu^*:\operatorname{Pic}_{X/\Omega}\to
\operatorname{Pic}_{Z/\Omega}$ is annihilated by $p^e$ as an fppf
sheaf.  Indeed, after an fppf localization on a test scheme $T$, an
element of this kernel is represented by a line bundle $L$ on $X_T$
such that $\nu_T^*L$ is pulled back from $T$.  The two semilinear maps
$\sigma$ and $F_T^e$ define a morphism
$\sigma\times F_T^e:X_T\to Z_T$.  Their composite with $\nu_T$ is the
absolute Frobenius $F_{X_T}^e$.  Pulling back $\nu_T^*L$ along this
morphism shows that $L^{\otimes p^e}$ is pulled back from $T$.
This proves the claim uniformly for all $T$; no identification of a
relative Frobenius with multiplication on a Picard scheme is used.

The normal proper $\Omega$-scheme $Z$ has proper torsion Picard
component, by Kleiman's Theorem~5.4 and Remarks~5.6 and~6.19
\cite{Kleiman2005}.  A homomorphism from a torus to a proper algebraic
group is zero: its scheme-theoretic image is a torus quotient and is
also proper, hence trivial.  Therefore any torus in
$\operatorname{Pic}^{\tau}_{X/\Omega}$ lies in the kernel of $\nu^*$.
Multiplication by $p^e$ is surjective on a torus, whereas it is zero on
this kernel.  The torus is consequently trivial.  Picard base change
identifies the group in (i) with $\operatorname{Pic}^{\tau}_{X/\Omega}$.

For (ii), take the closed scheme-theoretic images
\[
 H_n=\operatorname{im}([p^n]:H^0\longrightarrow H^0).
\]
The homomorphism theorem makes each map onto its image faithfully
flat \cite[Theorem 5.39, p.~108]{MilneAG2022B}.  The $H_n$ form a
descending chain of closed subgroup schemes in the noetherian scheme
$H^0$, so their defining ideal sheaves form an ascending chain and
stabilize.  Write $K=H_N=H_{N+1}=\cdots$.  Thus
$[p]:K\to K$ is faithfully flat.  Scheme-theoretic images commute with
flat field extension, so these assertions remain true over $\Omega$.

Over $\Omega$, the reduction $K_{\Omega,\mathrm{red}}$ is a smooth
connected subgroup.  The quotient
$Q=K_\Omega/K_{\Omega,\mathrm{red}}$ is finite and infinitesimal.
It is annihilated by a power of $p$.  For completeness, its local
augmentation ideal $I$ is nilpotent, and $[p]^*I\subset I^2$, since
$d[p]=p\,\mathrm{id}=0$ on the cotangent space.  Iteration sends $I$
into $I^{2^r}$ and is zero for large $r$.
But $[p]$ is faithfully flat on $Q$, so a power of it can be zero only
when $Q=0$.  Hence $K_\Omega$ is smooth.

Apply the Barsotti--Chevalley theorem over $\Omega$
\cite[Theorem 8.27, pp.~154--155]{MilneAG2022B}:
\[
 0\longrightarrow U\longrightarrow K_\Omega\longrightarrow B
 \longrightarrow0,
\]
where $U$ is smooth connected affine and $B$ is an abelian variety.
The group $U$ has no torus, so it is unipotent.  Some $[p^a]$ kills
$U$: embed it in a group of upper triangular unipotent matrices and
choose $p^a$ at least the matrix size.  Therefore $[p^a]$ on $K_\Omega$
factors through $B$.  It is surjective, so $K_\Omega$ is proper.
It follows that $K$ is a smooth proper connected $k$-group, hence an
abelian variety.  This deduction is made on the subgroup $K$ already
defined over $k$; it does not assume that a maximal abelian subvariety
commutes with an imperfect field extension.

Any abelian subvariety $A\subset H$ lies in $H^0$ and satisfies
$[p^N]A=A$, whence $A\subset K$.  Thus $K=A_H$ is the largest such
subvariety.  Finally $H/H^0$ is a finite etale group scheme.  If $c$
annihilates it, then $[cp^N]H\subset A_H$, proving the bounded-exponent
assertion including all infinitesimal structure and components.

For (iii), restriction of scalars exists as a finite-type algebraic
group: a proper algebraic group is quasi-projective, and restriction
along a finite field extension is representable for quasi-projective
schemes.  This is the restriction-of-scalars construction reviewed in
\cite[Section 2i, pp.~56--59]{MilneAG2022B}.
Over $\Omega$, write
$l\otimes_k\Omega=\prod_\sigma A_\sigma$, where each $A_\sigma$ is
local Artinian with residue field $\Omega$.  Restriction of scalars
commutes with this product and with field extension.  For each factor,
write $J_\sigma$ for its nilpotent augmentation ideal; choose a common
$r$ with $J_\sigma^r=0$ for all $\sigma$.  Reduction gives a homomorphism
\[
 \epsilon:(\operatorname{Res}_{l/k}P)_\Omega
       \longrightarrow\prod_\sigma P\otimes_{l,\sigma}\Omega.
\]
We show that $\ker\epsilon$ is killed by a fixed power of $p$.
We work on one factor and write $A=A_\sigma$ and $J=J_\sigma$.
For any $\Omega$-algebra $R$, kernel points on this factor are maps
$\operatorname{Spec}(R\otimes_\Omega A)\to P_A$ reducing to the identity
modulo $R\otimes_\Omega J$.  Such a map has image in an affine
neighborhood of the identity.  On the completed local algebra along
the identity section, $[p]^*$ sends the augmentation ideal into its
square, again because $d[p]=0$.  Therefore $[p]^e$ sends an identity
congruence modulo $R\otimes J$ to one modulo $R\otimes J^{2^e}$.
Choose $2^e\ge r$.  The same $[p^e]$ then kills the kernel on every
test algebra $R$, hence as a group scheme.  This argument does not
require $P$ to be smooth or the evaluation map to be surjective.

The same exponent works on all finitely many factors.  Any geometric
torus in the restriction of scalars maps trivially to the proper
product $\prod_\sigma P\otimes_{l,\sigma}\Omega$.  It is therefore contained in this
bounded-$p$-exponent kernel, and so is trivial.  A closed subgroup
inherits the same exclusion, completing the proof.
\end{proof}

\begin{remark}
This proposition is an auxiliary result for the weak Picard
stabilization argument.  It does not by itself prove that images of
Picard schemes on the permitted model category stabilize.  Its
bounded-exponent assertion is a statement about group schemes, not
only their geometric points.
\end{remark}


% Independently audited complete proof; amendments make descent points explicit.
\subsubsection*{Statement B.2s (Proved): weak stabilization when a regular birational model exists}
\begin{theorem}
\label{B:B2-regular-model-stabilization}
Let $k$ be a field of characteristic $p>0$, and let $K/k$ be a
finitely generated regular extension.  Suppose that $K/k$ has a
proper regular model $Y$.  Given any proper geometrically normal
model $X_*$ of $K/k$, there is a proper geometrically normal
modification $X_0\to X_*$ such that every proper geometrically
normal modification $W\to X_0$ induces an isomorphism
\[
 \operatorname{Pic}^{\tau}_{X_0/k}
       \xrightarrow{\ \sim\ }\operatorname{Pic}^{\tau}_{W/k}.
\]
The assertion concerns an upper cone and does not assert that this
cone is cofinal in the entire model category.
\end{theorem}

\begin{proof}
Write $P=\operatorname{Pic}^{\tau}_{Y/k}$.  This is a commutative
group scheme of finite type; it need not be proper over $k$.
For every permitted model $X\to X_*$, normalize the main graph
between $X$ and $Y$, obtaining $T\to X,Y$.  Statement B.1g
identifies $\operatorname{Pic}^{\tau}_T$ with $P$.
Since the proper birational map $T\to X$ has universal direct
image of its structure sheaf equal to $\mathcal O_X$, pullback
gives a closed immersion
\[
                       P_X:=\operatorname{Pic}^{\tau}_{X/k}
                              \lhook\joinrel\longrightarrow P.
\]
Here $P_X$ is proper because $X$ is geometrically normal.
A further normal common graph shows that the immersion is
independent of the choice of $T$, and that these immersions are
compatible with domination.  We henceforth regard the $P_X$ as
closed subgroup schemes of $P$.

Choose $X_1\to X_*$ for which $\dim P_{X_1}$ is maximal.
Such a choice exists because these dimensions are nonnegative
integers bounded by $\dim P$.  Let $A_0$ be the largest abelian
subvariety of $P_{X_1}$, obtained from the stable multiplication
image in B.2r.  Since $P_{X_1}$ is proper, its geometric reduced
identity component is an abelian variety; multiplication by $p$
is surjective on that component.  The stable image therefore
has dimension $\dim P_{X_1}$, so $P_{X_1}/A_0$ is finite.
Every subsequent $X\to X_1$ has the same largest abelian
subvariety $A_0$ inside $P$: the inclusion of these abelian
subvarieties has equal dimensions, again using properness.
Statement B.2r shows that $P$ has no geometric torus.

We construct, after finitely many further permitted modifications,
a closed subgroup $H\subset P$ containing all the subsequent
$P_X$ and having maximal abelian subvariety exactly $A_0$.
Initially take $H=P$.  The largest abelian subvariety $B\subset H$
exists by the stable-multiplication-image construction in B.2r.
It contains $A_0$.  If it equals $A_0$, the construction is finished.
Suppose instead that $\dim B>\dim A_0$.

\par\smallskip\noindent\textbf{A bad divisor attached to the excessive abelian part.}\quad
The Albanese torsor of the proper geometrically integral variety
$Y$ is dual to the maximal abelian subvariety of its Picard
scheme \cite[Theorem 3.2 and Remark 3.3, pp.~12--13 of the
author manuscript]{C_ACMV2026}.  Pushing this torsor out along
the dual of the inclusion of $B$ gives a torsor $\mathcal A_B$
under $B^\vee$ and a morphism $Y\to\mathcal A_B$ whose pullback
on $\operatorname{Pic}^0$ is the inclusion $B\subset P$.
Thus there is a rational map $X_1\dashrightarrow\mathcal A_B$.
Using the torsor avoids any extension of the ground field to
choose an origin.

Let $G$ be the normalization of its closed graph.  This is a
proper normal model with maps to $X_1$ and $\mathcal A_B$.
If $G$ were geometrically normal, it would be a permitted
modification of $X_1$ and its Picard image would contain $B$,
contrary to the maximality of $\dim A_0$.
Consequently $G$ is not geometrically normal.
Normality implies geometric $S_2$ under field extensions
\cite[Tags 045M and 0339]{StacksB}.  Geometric integrality
follows from the regularity of $K/k$.  By Serre's criterion,
there is a codimension-one point of $G$ whose local ring fails
geometric normality.  More explicitly, the failure occurs after
a finite algebraic extension, and a bad codimension-one point
there lies above a codimension-one point of $G$.
Denote the resulting divisorial valuation of $K$ by $v$ and its
residue field by $\kappa(v)$.

This valuation cannot have a divisorial centre on any permitted
model.  Otherwise the common normal graph over that model and
$G$ would be an isomorphism near the two codimension-one DVRs,
identifying the bad local ring with a geometrically normal one.
Its centre therefore has dimension at most $n-2$, where
$n=\operatorname{trdeg}_k K$.
The divisor on $G$, on the other hand, has image of dimension
$n-1$ in $X_1\times\mathcal A_B$, because normalization of
the graph is finite over its image.

Among permitted modifications of $X_1$, choose $X_2$ on which
the centre $C_2$ of $v$ has maximal dimension.  For every
permitted $W\to X_2$, its centre $C_W$ maps dominantly and
properly to $C_2$ and has the same dimension.  In particular
\[
                 F:=k(C_2)\ \subset\ k(C_W)
\]
is a finite extension inside the fixed residue field $\kappa(v)$.
We use the reduced centre in this assertion.

\par\smallskip\noindent\textbf{Restriction to a fixed proper generic fibre.}\quad
Take a normal proper common graph $T$ dominating $G,Y,X_2$.
Let $D$ be the normalization of the strict transform on $T$
of the divisor defining $v$.  It is proper over the centre
$C_2$, and the generic fibre $D_F$ is a proper normal integral
$F$-scheme with function field $\kappa(v)$.
Pullback from $Y$ followed by restriction defines a homomorphism
\[
                P_F\longrightarrow\operatorname{Pic}_{D_F/F}.
\]
Let $K_F$ be its closed scheme-theoretic kernel.

There is a largest closed $k$-subgroup $K_v\subset P$ whose
base change to $F$ factors through $K_F$, with the corresponding
factorization property on every test scheme.
To see this directly, on an affine open $\operatorname{Spec}R$
of $P$ choose finitely many generators of the ideal of $K_F$
in $R\otimes_k F$.  Expand the finitely many coefficients in
a $k$-basis of $F$ and let $I\subset R$ be generated by their
coefficients in $R$.  For every $k$-algebra $S$, a map $R\to S$
annihilates $I$ exactly when its base change to $F$ annihilates
the original ideal.  Thus these closed subschemes glue; the
group law, inverse and identity restrict by this universal
property.  No Weil restriction along an infinite field
extension is being asserted.

Every $P_W$ for a permitted $W\to X_2$ is contained in $K_v$.
Indeed, take a further normal common graph over $T$ and $W$,
and normalize the strict transform $D'$ of $D$.  The morphism
$D'_F\to D_F$ is proper birational between normal integral
schemes.  Its direct image of the structure sheaf
is $\mathcal O_{D_F}$; this remains so on every parameter base
change, since every $F$-scheme is flat over the field $F$.
Thus pullback of relative Picard sheaves
is injective.  The map $D'_F\to W$ factors through
$\operatorname{Spec}k(C_W)$ at the generic fibre of $C_2$.
A line bundle pulled back from $W$ is therefore, on this
generic fibre, pulled back from this finite field extension of
$F$.  Its relative Picard class is zero on all parameter
schemes: a line bundle on a finite locally free algebra becomes
trivial after the smooth surjective generator-open cover used
in B.2m.  This argument includes nonreduced parameters.
Injectivity of pullback to $D'_F$ proves the assertion on $D_F$,
and the defining property of $K_v$ proves $P_W\subset K_v$.

The subgroup $B$ is not contained in $K_v$.
Indeed the image of $D$ in $X_1\times\mathcal A_B$ still has
dimension $n-1$, whereas $\dim C_2\le n-2$.
The image of $D$ in $C_2\times\mathcal A_B$ dominates this
first image and therefore has dimension at least $n-1$.
Consequently the restriction $D_F\to(\mathcal A_B)_F$ has
positive-dimensional image.  After an algebraic closure of
$F$, at least one reduced irreducible component has a
nonconstant map to the torsor, and so does its normalization.
It is a proper geometrically integral variety.  The Albanese
universal property and Picard duality show that a nonconstant
map from it to an abelian torsor has nonzero pullback on
$\operatorname{Pic}^0$.  Since this pullback is the restriction
of $B\to P_F$, we obtain $B\not\subset K_v$.

Replace $H$ by $H\cap K_v$ and the current model by $X_2$.
At the next step we relabel this current model as $X_1$;
the subgroup $A_0$ stays fixed throughout the construction.
All later Picard images lie in the new $H$.  Its largest
abelian subvariety is contained in the proper closed subgroup
$B\cap K_v$ of $B$, and hence has strictly smaller dimension.
It still contains $A_0$.  Repeating this construction can take
at most $\dim B-\dim A_0$ steps.  We therefore reach a permitted
model $X_3$ and a fixed closed subgroup $H\subset P$ such that
all $P_W$, for $W\to X_3$, lie in $H$ and its largest abelian
subvariety is $A_0$.

\par\smallskip\noindent\textbf{The remaining finite quotient has a uniform order bound.}\quad
By B.2r there is an integer $N>0$ with $[N]H\subset A_0$.
For every later $W$, the quotient $P_W/A_0$ is finite, and
the map $P_W[N]\to P_W/A_0$ is fppf-surjective.
For an explicit verification, if $x$ is a local section of
$P_W$, then $Nx\in A_0$; fppf locally we can choose
$a\in A_0$ with $Na=Nx$, since multiplication by $N$ on an
abelian variety is faithfully flat.  The section $x-a$ belongs
to $P_W[N]$ and represents the same quotient class.

After one algebraically closed extension, every such $P_W$
embeds compatibly into a fixed $g$-dimensional Jacobian $J$ by
B.2d.  To apply that statement, one may use normal graphs over
one fixed projective model after algebraic closure; no
geometrically normal common refinement over $k$ is needed.
We obtain
\[
       |P_W/A_0|\le |P_W[N]|\le |J[N]|=N^{2g}.
\]
Choose $X_0\to X_3$ for which the finite order $|P_{X_0}/A_0|$
is maximal.  Every later closed immersion
$P_{X_0}\hookrightarrow P_W$ induces a closed immersion of
finite quotients of the same order, hence an isomorphism.
It follows that $P_{X_0}\to P_W$ itself is an isomorphism.
This proves the theorem without a filteredness assumption.
\end{proof}

\begin{corollary}
The weak upper-cone assertion holds over every field of positive
characteristic for function fields of transcendence degree at most
three, whenever the class of proper geometrically normal models
is nonempty.
\end{corollary}
\begin{proof}
For a geometrically integral function field in dimension at most
three, a proper regular birational model exists by
\cite[Theorem 1.1, pp.~3--4 of the author manuscript]{CossartPiltant2019B}.
The hypotheses of that theorem include finite-type schemes over
arbitrary fields: they are excellent, and the asserted resolution
is proper and birational.
If $K$ has nontrivial algebraic constants $k_0/k$, geometric
normality of a permitted model implies that $k_0/k$ is finite
separable.  We apply the result over $k_0$, where $K/k_0$ is
regular, and descend the resulting Picard isomorphisms through
$\operatorname{Res}_{k_0/k}$, as in B.2o.
\end{proof}


\subsubsection*{Statement B.2t (Proved): abelian realization through a Galois alteration}

\begin{lemma}[Equivariant alteration data]
\label{B:equivariant-constant-data}
Let $k$ have characteristic $p>0$, let $K/k$ be regular and
finitely generated, and let $X_0$ be a proper geometrically
normal model.  We may choose a finite field extension $l/k$,
a smooth projective geometrically integral $l$-scheme $Y$, a
proper generically finite $k$-morphism $Y\to X_0$, and a finite
group $\Gamma$ acting on this morphism, with the following
properties.  Writing $L=k(Y)$, the field $l$ is stable under
$\Gamma$, its action on $Y$ is semilinear over this action on
$l$, and
\[
 K'=L^\Gamma\quad\hbox{is finite purely inseparable over }K.
\]
In particular $l_0=l^\Gamma$ is finite purely inseparable over
$k$, and $l/l_0$ is finite Galois.  The group action on $l$
need not be trivial.
\end{lemma}
\begin{proof}
Apply the regular Galois alteration theorem over $k^{\mathrm{perf}}$
to $X_0\times_k k^{\mathrm{perf}}$; use an initial projective
birational modification if needed.  The precise primary
statements are Situation~5.3, Theorem~5.13 and Corollary~5.15 of
\cite[pp.~613 and 619--620]{deJong1997}; projectivity is retained
in Theorem~5.9 and Proposition~5.11.  We obtain a regular
projective integral $\widetilde Y$ and a finite group $\Gamma$
with invariant function field purely inseparable over
$Kk^{\mathrm{perf}}$.  As the base field is perfect,
$\widetilde Y$ is smooth.  Let $\widetilde l$ be its field of
algebraic constants.  It is a finite separable extension of
$k^{\mathrm{perf}}$, and $\widetilde Y$ is geometrically
integral over $\widetilde l$.

The group preserves $\widetilde l$ and
$\widetilde l^\Gamma=k^{\mathrm{perf}}$.  Indeed an algebraic
constant in the invariant field is purely inseparable over
$Kk^{\mathrm{perf}}$ and separable over the perfect constant
field.  Regularity of $Kk^{\mathrm{perf}}/k^{\mathrm{perf}}$
forces this constant to lie in $k^{\mathrm{perf}}$.

Choose a finite extension $l/k$ inside $\widetilde l$ containing
a primitive generator of $\widetilde l/k^{\mathrm{perf}}$ and
all coefficients needed for the scheme, the map, the group
action and its finitely many action identities.  Adjoin their
finitely many $\Gamma$-conjugates.  Enlarging this finite field
if necessary, finite-presentation descent gives all the
specified data over $l$, and
$\widetilde l=lk^{\mathrm{perf}}$.  Smoothness, projectivity and
geometric integrality of $Y/l$ follow by faithfully flat field
descent.  Its map to $X_0$ is dominant of the same dimension,
hence generically finite, and is proper.  Thus $L/K$ is finite.

There is a compatible embedding $L\subset k(\widetilde Y)$.
Consequently $L^\Gamma$ lies in the purely inseparable extension
$k(\widetilde Y)^\Gamma$ of $Kk^{\mathrm{perf}}$.
Since $Kk^{\mathrm{perf}}/K$ is purely inseparable,
$L^\Gamma/K$ is purely inseparable; it is finite since $L/K$
is finite.  The same argument gives
$l^\Gamma=l\cap k^{\mathrm{perf}}$.  The final Galois assertion
is the fixed-field theorem for a finite group of field
automorphisms.  This construction expressly allows the finite
separable constants that an alteration over a perfect field
can introduce.
\end{proof}

\begin{lemma}[Realization of invariant abelian subgroups]
\label{B:invariant-abelian-realization}
Choose the preceding data and put
\[
 P=\operatorname{Pic}^{\tau}_{Y/l},\qquad
 \mathcal R=\operatorname{Res}_{l/k}(P),\qquad
 \mathcal A=\mathcal R^\Gamma.
\]
Let $H\hookrightarrow\mathcal A$ be an abelian variety over
$k$.  There is a torsor $Q$ over $k$ under an abelian variety,
and a rational map $X_0\dashrightarrow Q$, such that its
pullback on connected Picard schemes, evaluated on the fixed
alteration, has scheme-theoretic image exactly $H$ in
$\mathcal R$.

More precisely, let $\Gamma_0$ be the kernel of the action on
$l$, put $n=|\Gamma_0|$, and choose $q=p^e$ with $K'^q\subset K$.
We may take $Q$ to be a torsor under $(H^\vee)^{(q)}$, and the
resulting homomorphism is
\[
 H^{(q)}\xrightarrow{V_H^e}H\xrightarrow{[n]}H
       \hookrightarrow\mathcal R.                 \tag{B.2.14}
\]
Here $V_H^e$ is the dual of the $e$-fold relative Frobenius of
$H^\vee$; it is an isogeny.  In particular the image assertion
does not require either $n$ or $q$ to be prime to $p$.
\end{lemma}
\begin{proof}
\par\smallskip\noindent\textbf{Evaluation has finite kernel on $H$.}\quad
Adjunction gives $u:H_l\to P$.  We check that its kernel is
finite.  Over an algebraic closure $\Omega$ of $k$, write
\[
 l\otimes_k\Omega=\prod_{\sigma}A_\sigma,
\]
where each $A_\sigma$ is an Artin local $\Omega$-algebra with
residue field $\Omega$.  Reduction defines a map from
$\mathcal R_\Omega$ to the product of the corresponding
geometric Picard schemes.  Its kernel is affine.  To see this
without a smoothness assertion about $P$, choose an affine
neighbourhood of the identity in each base-changed Picard
scheme over $A_\sigma$.  Every morphism from $T_{A_\sigma}$
whose reduction is the identity lands in that neighbourhood,
because $T\to T_{A_\sigma}$ is a nilpotent closed immersion
and a homeomorphism.  The kernel is therefore a closed
subscheme of the Weil restriction of an affine scheme along
a finite locally free map, which is affine.

The group $\Gamma$ acts transitively on the embeddings indexed
by $\sigma$, since $l^\Gamma/k$ is purely inseparable.
As $H$ lies in the pointwise fixed subgroup, the kernels of all
its geometric evaluation maps are equal.  Their common kernel
is contained in the affine reduction kernel just described.
It is also a closed subgroup of the proper scheme $H_\Omega$.
It is therefore proper and affine, and hence finite.
This proves the assertion for $u$ by field descent.

\par\smallskip\noindent\textbf{The Albanese torsor and its semilinear action.}\quad
The image of $u$ is an abelian subvariety of $P$.
Duality for the proper geometrically integral scheme $Y/l$
gives a surjective homomorphism
$\operatorname{Alb}_{Y/l}\to H_l^\vee$.
Push out the Albanese torsor along this homomorphism.  We get
an $H_l^\vee$-torsor $T/l$ and a morphism $f:Y\to T$ whose
pullback on connected Picard schemes is $u$.
We use the proper geometrically reduced Albanese duality in
\cite[Theorem~3.2 and Remark~3.3, author pp.~12--13]{C_ACMV2026}.

The invariance $H\subset\mathcal R^\Gamma$ and functoriality
of the Albanese torsor give a semilinear $\Gamma$-action on
$T$ making $f$ equivariant.  Its linear action on the acting
abelian variety is the natural semilinear action on
$H^\vee_l$, because $\Gamma$ fixes the $k$-group $H$ pointwise.
In particular every $\gamma\in\Gamma_0$ acts on $T$ by a
translation $c_\gamma\in H^\vee(l)$.  The assignment
$\gamma\mapsto c_\gamma$ is a homomorphism: $\Gamma_0$ acts
trivially on $l$ and on $H^\vee$.  Thus $[n]c_\gamma=0$.

Push $T$ out along $[n]:H_l^\vee\to H_l^\vee$, obtaining
$T_n$ and $f_n:Y\to T_n$.  The subgroup $\Gamma_0$ acts
trivially on $T_n$.  The remaining semilinear action is the
Galois descent datum for $l/l_0$, where $l_0=l^\Gamma$.
It descends $T_n$ to an $H^\vee_{l_0}$-torsor $T_0/l_0$.
On function fields the equivariant map $f_n$ gives a
$\Gamma$-fixed point of $T_0(L)$, hence a point of $T_0(K')$.

\par\smallskip\noindent\textbf{Purely inseparable descent by a field homomorphism.}\quad
The inclusion $l_0\subset K'$ and $K'^q\subset K$ imply
$l_0^q\subset k$: the field $k$ is algebraically closed in
$K$.  Hence $a\mapsto a^q$ is a field homomorphism
$l_0\to k$.  Define
\[
 Q=T_0\times_{\operatorname{Spec}l_0, a\mapsto a^q}
                   \operatorname{Spec}k.
\]
It is a torsor under $(H^\vee)^{(q)}$.
Base-changing the $K'$-point of $T_0$ along the field
homomorphism $K'\to K$, $b\mapsto b^q$, gives a $K$-point
of $Q$, and therefore a rational map $X_0\dashrightarrow Q$.

The canonical $k$-morphism $T_0\to Q$ has, on coordinate
algebras, the formula $b\otimes c\mapsto b^q c$.
After extension to $l_0$ it is the relative Frobenius on the
torsor.  Its pullback on connected Picard schemes is
$V_H^e:H^{(q)}\to H$.  Combining it with the pushout $[n]$
and the original adjunction homomorphism $u$ gives precisely
(B.2.14).  Both $V_H^e$ and $[n]$ are finite faithfully flat
isogenies.  Their composite has image all of $H$, proving
the claimed scheme-theoretic image in $\mathcal R$.
\end{proof}

\begin{lemma}[Relative Picard sheaves with finite constants]
\label{B:Pic-finite-constants-sheaf}
Let $l/k$ be finite and let $Z/l$ be proper and geometrically
integral, with a nonempty smooth open.  There is a canonical
isomorphism of full fppf Picard sheaves
\[
 \operatorname{Pic}_{Z/k}\simeq
 \operatorname{Res}_{l/k}\operatorname{Pic}_{Z/l}.
\]
No identification of their torsion components is asserted.
\end{lemma}
\begin{proof}
Choose a nonempty affine smooth open $U\subset Z$.
The affine scheme $S=\operatorname{Res}_{l/k}U$ is smooth
and surjective over $k$.  Indeed finite locally free restriction
preserves formal smoothness and finite presentation; over an
algebraic closure the algebra $l\otimes_k\bar k$ is a product
of Artin local algebras, and a point of each smooth nonempty
geometric fibre lifts through its nilpotent thickening.
The universal point gives a section of $Z_{S_l}\to S_l$.

After the fppf base change $T\times_kS\to T$, every class in
$\operatorname{Pic}_{Z/l}(T_l)$ is represented by a line bundle
on $Z_{T_l}$.  For clarity, the section permits rigidification:
the universal equality $H^0(Z,\mathcal O_Z)=l$ makes the
automorphism of a rigidified line bundle trivial, so the local
representatives and their rigidifications descend effectively.
Two representatives differ by a line bundle from $T_l$.
Such a line bundle becomes trivial after an fppf cover of $T$,
by the smooth-surjective generator-open argument of B.2m.
Thus sending a representative to its $k$-relative Picard class
is well-defined as a map of fppf sheaves and is inverse to the
natural map in the other direction.
\end{proof}

\begin{remark}[Application to the fixed comparison group]
Every comparison map $h_X$ of B.2m lands in
$\mathcal R^\Gamma$: the normal main graph $Y_X$ carries the
group action, its comparison with the regular $Y$ is
equivariant, and its pullback from $X$ is invariant.
Moreover the preceding full-sheaf identification applies to
every normal main alteration graph $Y_X/l$: its function field
is regular over $l$, so it has a nonempty smooth open.  Combining
it with B.1g over $l$ maps the entire comparison group
$\mathcal R$ into $\operatorname{Pic}_{Y_X/k}$, making the
generic-fibre restriction kernels used in the stabilization
argument defined on the full ambient group.
The construction retains the possible separable constants of an
alteration over $k^{\mathrm{perf}}$.  Its semilinear action,
translation quotient and purely inseparable field homomorphism
are part of the proof, rather than an assumption that those
constants disappear.  All three lemmas have undergone the
targeted proof checks recorded in the accompanying ledger.
\end{remark}


% Complete proof after independent actual-file audits.
\subsubsection*{Statement B.2u (Proved): weak upper-cone stabilization over every field}
\begin{theorem}
\label{B:B2-arbitrary-field-stabilization}
Let $k$ have characteristic $p>0$, let $K/k$ be finitely generated,
and suppose its class of proper geometrically normal models is
nonempty.  Given any such model $X_*$, there is a proper
geometrically normal modification $X_0\to X_*$ such that
\[
 \operatorname{Pic}^{\tau}_{X_0/k}\xrightarrow{\ \sim\ }
 \operatorname{Pic}^{\tau}_{W/k}
\]
for every proper geometrically normal modification $W\to X_0$.
This is the weak upper-cone assertion.  A constant domination-cofinal
diagram can still fail, as B.2q proves.
\end{theorem}

\begin{proof}
We first assume $K/k$ regular.  Choose the equivariant alteration
data of Lemma~\ref{B:equivariant-constant-data} over $X_*$:
a smooth projective geometrically integral $Y/l$, with $l/k$
finite, function field $L$, and a finite semilinear group
$\Gamma$ such that $L^\Gamma/K$ is finite purely inseparable.
Put
\[
 d=[L:K],\qquad P=\operatorname{Pic}^{\tau}_{Y/l},\qquad
 R=\operatorname{Res}_{l/k}P,\qquad A=R^\Gamma.
\]
For each permitted $X\to X_*$, normalize the main graph
$Y_X$ between $Y$ and $X$.  It has function field $L$, and
$Y_X\to Y$ is proper birational with normal source.  By B.1g
over $l$ its torsion Picard group is canonically $P$.
This identification respects the semilinear $\Gamma$-action.
Consequently B.2m gives compatible homomorphisms
\[
 h_X:P_X:=\operatorname{Pic}^{\tau}_{X/k}\longrightarrow A
           \subset R
\]
whose kernels are finite and killed by $d$.  The assertion
about finite kernels includes all parameter schemes, by the
fixed-Jacobian bound in B.2m.  Pullbacks among the $P_X$ are
closed immersions.  The group $R$ is of finite type and has
no geometric torus by B.2r(iii), which allows arbitrary finite
constant extensions.

Choose $X_1\to X_*$ maximizing $\dim P_X$.
The dimensions are bounded by $\dim R$, because the kernels
of $h_X$ are finite.  Let $A_*\subset P_{X_1}$ be its largest
abelian subvariety, furnished by B.2r.  Properness of $P_{X_1}$
implies that $P_{X_1}/A_*$ is finite.  Every subsequent
$W\to X_1$ has this same abelian subvariety under pullback,
since its dimension is maximal; we identify those copies of
$A_*$ without identifying the entire $P_W$.
Let
\[
                         A_0=h_{X_1}(A_*)\subset A.
\]
This is an abelian variety isogenous to $A_*$, since the
kernel is finite.  The distinction between $A_*$ and its
image $A_0$ is needed below.

\par\smallskip\noindent\textbf{Removing excessive abelian directions in the ambient group.}\quad
Initially set $H=A$.  We will impose finitely many closed
subgroup conditions on $H$ and replace the current permitted
model by a domination, always keeping every subsequent image
$h_W(P_W)$ inside $H$.  Let $B\subset H$ be its largest
abelian subvariety, which exists by B.2r and contains $A_0$.
If $B=A_0$, no further condition is needed.
Otherwise $\dim B>\dim A_0$.

Lemma~\ref{B:invariant-abelian-realization} supplies an
abelian torsor $Q/k$ and a rational map $X_1\dashrightarrow Q$
whose pullback on connected Picard groups, evaluated in $R$,
has scheme-theoretic image exactly $B$.
Let $G$ be the normalization of the closed graph.
If it were geometrically normal, it would be a permitted
modification and its image $h_G(P_G)$ would contain $B$.
This contradicts maximality of $\dim P_{X_1}=\dim A_0$.
Thus $G$ is not geometrically normal.

Exactly as in B.2s, normality gives geometric $S_2$, and
regularity of $K/k$ gives geometric integrality.  Hence there
is a geometrically bad divisorial valuation $v$ on $G$.
Its centre on every permitted model has dimension at most
$n-2$, where $n=\operatorname{trdeg}_k K$, because a proper
birational normal morphism is an isomorphism near a
codimension-one point.  The divisor on $G$ has image of
dimension $n-1$ in $X_1\times Q$.
Choose a permitted $X_2\to X_1$ on which the centre $C_2$
of $v$ has maximal dimension.  For every later $W\to X_2$,
its centre $C_W$ has that same dimension, and
\[
                     F=k(C_2)\ \subset\ k(C_W)
\]
is a finite extension inside $\kappa(v)$.

Choose an extension $w$ of $v$ to the finite field extension
$L/K$.  It is a divisorial valuation of $L$.
Normalize $G$ in $L$ and take a normal proper common graph
$Z$ over this normalization, $Y$ and $X_2$, with function
field $L$.  By Chow's lemma followed by normalization, we
may and do take $Z$ projective over $k$
\cite[Tag 0200]{StacksB}; no geometric normality of $Z$ is
required.  The map $Z\to Y$ is birational, and the divisor
of $w$ persists because the proper birational map to the
normalization of $G$ is an isomorphism at its generic point.
Let $D$ be the normalization of that divisor on $Z$.
Then $D_F$ is proper normal integral over $F$, with function
field $\kappa(w)$, finite over $\kappa(v)$.
It is projective, so the relative Picard group and its closed
scheme-theoretic kernels have the usual projective
representability; the same conclusion also follows from
proper Picard representability over a field.

The restriction condition must be defined on all of $R$,
rather than merely on the known Picard images.  B.1g over
$l$ identifies $P$ with $\operatorname{Pic}^{\tau}_{Z/l}$.
The function field $L/l$ is regular, so $Z/l$ has a nonempty
smooth open.  Lemma~\ref{B:Pic-finite-constants-sheaf} gives
the full fppf identification
\[
 \operatorname{Pic}_{Z/k}
       \simeq\operatorname{Res}_{l/k}\operatorname{Pic}_{Z/l}.
\]
It follows that restriction gives a homomorphism
\[
                       R_F\longrightarrow\operatorname{Pic}_{D_F/F}.
\]
This uses no equality between torsion components under
restriction of scalars.  Let $K_w\subset R$ be the closed
$k$-subgroup defined by the vanishing of this restriction
after base change to $F$; its coefficient-ideal construction
and universal factorization property are exactly those
proved in B.2s.

All later images $h_W(P_W)$ are contained in $K_w$.
To check this, form a further normal common graph over $Z$
and $W$ and normalize the strict transform $D'$ of $D$.
The map $D'_F\to D_F$ is proper birational between normal
integral schemes, so its direct image of the structure sheaf
is $\mathcal O_{D_F}$.  Every parameter base change over the
field $F$ is flat, making relative Picard pullback injective
on all test schemes.  The map from $D'_F$ to $W$ factors
through $\operatorname{Spec}k(C_W)$.
Thus a line bundle from $W$ restricts to one pulled back
from the finite $F$-algebra $k(C_W)$.
Its relative Picard class is zero after the smooth
surjective generator-open cover of the parameter scheme.
This includes nonreduced parameters and inseparable centre
field extensions.  Injectivity on $D'_F\to D_F$ proves the
required vanishing on the fixed fibre.

The restriction does not vanish on all of $B$.
The image of $D$ in $X_1\times Q$ has dimension $n-1$:
the extension $\kappa(w)/\kappa(v)$ is finite and does not
alter the dimension of the original divisor image.
The image in $C_2\times Q$ dominates that image, while
$\dim C_2\le n-2$.
Therefore $D_F\to Q_F$ has positive-dimensional image.
After algebraic closure and normalization of a reduced
irreducible component, Albanese duality shows that its
pullback on $\operatorname{Pic}^0_Q$ is nonzero.
If $B$ were contained in $K_w$, this pullback would be zero,
because its map to $R$ has image $B$.  This contradiction
proves $B\not\subset K_w$.

Replace $H$ by $H\cap K_w$ and relabel the current model
$X_2$ as $X_1$.  The largest abelian subvariety of the new
$H$ is contained in the proper subgroup $B\cap K_w$, so
its dimension strictly drops, while it still contains $A_0$.
After finitely many steps we obtain a permitted model $X_3$
and a closed subgroup $H\subset A$ containing every later
$h_W(P_W)$ and having largest abelian subvariety exactly $A_0$.

\par\smallskip\noindent\textbf{A degree factor controls the remaining comparison kernel.}\quad
By B.2r there is a fixed integer $N>0$ with
$[N]H\subset A_0$.
For every later $W$, we claim that
\[
                         [dN]P_W\subset A_*.
\]
Indeed, for a local section $x$ of $P_W$ the section
$N h_W(x)$ lies in $A_0$.
The isogeny $A_*\to A_0$ is fppf-surjective, so fppf locally
we can choose $a\in A_*$ with $h_W(a)=N h_W(x)$.
Then $Nx-a\in\ker h_W$, which is killed by $d$.
Consequently $dNx=da\in A_*$, proving the claim as a
group-scheme assertion.

The quotient $P_W/A_*$ is finite, since $A_*$ has the
maximal dimension of $P_W$.  It is killed by $dN$, and
\[
                  P_W[dN]\longrightarrow P_W/A_*
\]
is fppf-surjective: subtract a $(dN)$-division of
$(dN)x$ in the abelian variety $A_*$ from a local lift $x$.
After one fixed algebraically closed extension, B.2d embeds
all these $P_W$ into one $g$-dimensional Jacobian $J$;
normal graphs are required only over that algebraically
closed field.  Hence
\[
                    |P_W/A_*|\le |J[dN]|=(dN)^{2g}.
\]
Choose $X_0\to X_3$ with maximal finite order
$|P_{X_0}/A_*|$.  Every later closed immersion of Picard
groups induces a closed immersion between finite quotients
of that same order and is therefore an isomorphism.
This proves the theorem for regular $K/k$.

In general, the existence of a geometrically normal model
implies that the algebraic constant field $k_0$ of $K/k$
is finite separable over $k$, and that $K/k_0$ is regular.
We apply the theorem over $k_0$ to the model $X_*$.
Its underlying proper modifications are the same, and the
Picard isomorphisms over $k$ are obtained by
$\operatorname{Res}_{k_0/k}$, or checked after a finite
separable extension splitting the constants.  This completes
the arbitrary-field proof.
\end{proof}

\begin{remark}[The three formulations are distinct]
We have now proved both uniform $h^1$ boundedness and the
weak upper-cone assertion over arbitrary fields.  The
constant domination-cofinal formulation is false over an
imperfect field by B.2q, and the indicated model category
need not be filtered by B.2i.  The theorem does not replace
these counterexamples with a claim of cofinality: two
different stable upper cones may have nonisomorphic values.
\end{remark}
